\documentclass[11pt]{article}
\usepackage[margin=1in]{geometry}
\usepackage{amsmath,amssymb,amsthm,mathtools,booktabs,longtable,hyperref}
\newcommand{\Z}{\mathbb Z}
\newcommand{\F}{\mathcal F}
\newcommand{\Q}{\mathbb Q}
\newcommand{\Fp}{\mathbb F_p}
\newcommand{\vp}{v_p}
\newtheorem{theorem}{Theorem}
\newtheorem{lemma}{Lemma}
\newtheorem{proposition}{Proposition}
\newtheorem{corollary}[proposition]{Corollary}
\title{Hypergeometric readouts, a uniform common divisor,\\and a conditional irrationality criterion}
\author{}\date{}
\begin{document}\maketitle

\begin{abstract}
We give an explicit integer normalization for the harmonic readouts of
 a terminating hypergeometric trajectory. A family of factorial
permutations then supplies a common divisor valid simultaneously for
all integer linear combinations of the numerator readouts. The
construction is expressed directly in binomial coefficients and their
parameter derivatives. The supplied limit argument is separated into explicit integral,
matrix, and asymptotic statements. The final conclusion records the
precise dependencies that must be verified.
\end{abstract}

\paragraph{Purpose.}
This manuscript is a mathematical specification and proof draft for
formalization. It first explains the argument independently of any
trajectory, then gives the forward ${}_5F_4$ construction, and finally
supplies the analytic and arithmetic details. The attached scripts
provide exact algebraic certificates. No proof-assistant implementation
is part of this handoff.

\section{The conceptual argument}
Let $q_n>0,p_{3,n},p_{4,n}$ be rational numbers and set
\[
 r_n(A,B)=\frac{Ap_{3,n}+Bp_{4,n}}{q_n},\qquad
 \alpha_{A,B}=A\zeta(3)+B\zeta(4),\quad A,B\in\Z.
\]
The argument has three independent parts. Analysis identifies the limit
and bounds the error. Arithmetic constructs integers representing the
same rational approximants. A recurrence excludes an eventually zero
error for every nonzero pair $(A,B)$.

Choose a positive rational multiplier $m_n$ making
$m_n(q_n,p_{3,n},p_{4,n})$ an integer vector. Then find a positive integer
$g_n$ dividing its three coordinates. The resulting integers are
\[
 u_n=\frac{m_nq_n}{g_n}>0,\qquad
 v_n(A,B)=\frac{m_n(Ap_{3,n}+Bp_{4,n})}{g_n},\qquad
 r_n(A,B)=\frac{v_n(A,B)}{u_n}.
\]
This construction is simultaneous for all integer pairs: the multiplier
and common divisor are chosen for the three basic sequences, before
$A,B$ are selected. The integer $u_n$ need not be the reduced denominator;
the latter is at most $u_n$.

\begin{proposition}[The integer obstruction]\label{conceptual-criterion}
Fix a nonzero integer pair. Suppose
\[
 u_n\leq e^{\sigma n+o(n)},\qquad
 |\alpha_{A,B}-r_n(A,B)|\leq e^{-\beta n+o(n)},
 \qquad\beta>\sigma>0,
\]
and the error is nonzero for infinitely many $n$. Then
$\alpha_{A,B}$ is irrational. If the hypotheses hold for every fixed
nonzero pair, then $1,\zeta(3),\zeta(4)$ are linearly independent over
$\Q$. For reduced approximants, $\liminf\delta_n\geq\beta/\sigma-1$.
\end{proposition}
\begin{proof}
Along the nonzero errors, $u_n|\alpha_{A,B}-r_n|\to0$.
If $\alpha_{A,B}=a/b$ with integers $a,b$ and $b>0$, then
\[
 u_n|\alpha_{A,B}-r_n|
 =\frac{|au_n-bv_n|}{b}\geq\frac1b,
\]
a contradiction. Any rational relation among the three constants can
be cleared to an integer relation, giving the second assertion.
After irrationality is known all errors are nonzero and the reduced
denominators tend to infinity. Since $\log q_n^{\rm red}\leq\log u_n$,
the displayed exponential bounds imply the delta assertion.
\end{proof}

For a terminating hypergeometric sum, the coefficient multiplier $E_n$
is chosen to make every summand integral. Harmonic numbers appearing
in the numerator readouts require additional powers of an lcm $D_n$.
Permuting factorial factorizations gives different integer
normalizations of the same vector; comparing them prime by prime gives
a common divisor. An additional residue-sum identity improves that
divisor in this construction. Each of these roles is developed below.

\section{Coefficient normalization}
Let
\[
 \F_n(z)={}_pF_q(\mathbf a_n;\mathbf b_n;z)
       =\sum_{k=0}^{m_n}t_{n,k}z^k
\]
be a terminating hypergeometric series. Differentiation preserves its
hypergeometric form: for $0\leq s\leq m_n$,
\begin{equation}\label{differentiation}
 \frac{d^s}{dz^s}\F_n(z)
 =K_s(n)\F_{n,s}(z),\qquad
 K_s(n)=\frac{\prod_i(a_{n,i})_s}{\prod_j(b_{n,j})_s},\qquad
 \F_{n,s}(z)={}_pF_q(\mathbf a_n+s;\mathbf b_n+s;z).
\end{equation}
This is the standard differentiation formula
\cite[Eq.~16.3.1]{DLMF}. The index $s$ shifts the parameters, not the
trajectory depth $n$.

For $\theta=z\,d/dz$, the identity $\theta z^k=kz^k$ gives
\begin{equation}\label{polynomial-weight}
 [P(\theta)\F_n]_{z=1}=\sum_{k=0}^{m_n}P(k)t_{n,k}
 \qquad(P\in\Z[x]).
\end{equation}
Choose a positive integer $E_n$ that clears the coefficient denominators,
so that $B_{n,k}:=E_nt_{n,k}\in\Z$ for every $k$. Then
\begin{equation}\label{integer-weight}
 E_n[P(\theta)\F_n]_{z=1}=\sum_kP(k)B_{n,k}\in\Z.
\end{equation}
In the shifted representation \eqref{differentiation}, the same assertion
reads
\[
 E_nK_s(n)\F_{n,s}(1)
 =\sum_{k=s}^{m_n}\frac{k!}{(k-s)!}B_{n,k}\in\Z.
\]
Thus one coefficient normalization clears both the hypergeometric
series and the derivative prefactors.

Harmonic readouts introduce additional rational weights. Write
$H_m^{(r)}=\sum_{k=1}^m k^{-r}$ and
$d_M=\operatorname{lcm}(1,\ldots,M)$, with $d_0=1$.
We use the elementary integrality fact
\begin{equation}\label{harmonic-integrality}
 d_M^rH_m^{(r)}\in\Z\qquad(0\leq m\leq M).
\end{equation}
A common lcm will also clear the coefficient derivatives that occur
in these readouts. Alternative factorial factorizations of the same
coefficient family can give several such integer normalizations.
Comparing their prime valuations produces a common divisor; this is the
arithmetic permutation principle developed in the work of Rhin--Viola
and Zudilin \cite{RV,Zudilin}.

\section{The forward hypergeometric trajectory}
Use the coordinate order of the supplied convergence argument:
\[
 x_0=(0,0,1,1,1;1,1,1,1),\qquad
 v_F=(-14,-15,13,14,12;1,2,3,4).
\]
This is a simultaneous permutation of the upper parameters of the
earlier start and direction, so the defining hypergeometric polynomial
and the jet-basis contiguous equations agree. Set
\[
 \mathbf a_n=(-14n,-15n,1+13n,1+14n,1+12n),\qquad
 \mathbf b_n=(1+n,1+2n,1+3n,1+4n),
\]
\[
 \F_n(z)={}_5F_4(\mathbf a_n;\mathbf b_n;z)
 =\sum_{k=0}^{14n}
 \frac{\prod_i(a_{n,i})_k}{k!\prod_j(b_{n,j})_k}\,z^k.
\]
Let $T(n)$ be the native forward matrix obtained from the contiguous
relations and let $W_n=T(0)\cdots T(n-1)$, with $W_0=I$.
Its hypergeometric row has initial value $e_0^T$ and satisfies
\[
 e_0^TW_n=(\F_n(1),\theta\F_n(1),\theta^2\F_n(1),\theta^3\F_n(1)).
\]
In particular, the denominator row $r_Q=12e_0^T$ gives
$r_QW_ne_0=12\F_n(1)$.
For $A,B\in\Z$, the derivative observable
\[
 G_n(A,B)=[-6A\theta^2\F_n+(3A-2B)\theta^3\F_n]_{z=1}
\]
is $e_0^TW_nc_{A,B}$, where
$c_{A,B}=(0,0,-6A,3A-2B)^T$.
Choosing this vector as the first column of a basis matrix makes $G_n$
the first coordinate of a transformed row solution, with its initial
row transformed by the same matrix.
The original numerator row is $r_{A,B}=c_{A,B}^T$ and has readout
$r_{A,B}W_ne_0$; its harmonic representation is given below.

Formula \eqref{differentiation} gives the explicit shifted expression
\begin{align*}
 G_n(A,B)={}&(-3A-2B)K_1(n)\F_{n,1}(1)\\
 &+(3A-6B)K_2(n)\F_{n,2}(1)
 +(3A-2B)K_3(n)\F_{n,3}(1),
\end{align*}
where
\[
 K_1(n)=\frac{210n^2(1+12n)(1+13n)(1+14n)}
 {(1+n)(1+2n)(1+3n)(1+4n)}.
\]
For a quadratic weight, the equivalent single-series expression is
\[
 [(u\theta+v\theta^2)\F_n]_{z=1}
 =(u+v)K_1(n)\,{}_6F_5(\mathbf a_n+1,c+1;\mathbf b_n+1,c;1),
 \qquad c=(u+v)/v,
\]
whenever its parameters are regular. Polynomial weighting as in
\eqref{polynomial-weight} applies without these parameter restrictions.

For $n\geq1$, take
\begin{equation}\label{E}
 E_n=\binom{12n}{10n}\binom{13n}{10n}
        \binom{14n}{10n}\binom{15n}{n}.
\end{equation}
Factorial cancellation gives
\begin{equation}\label{B}
 B_{n,k}=E_nt_{n,k}
 =\binom{12n+k}{10n}\binom{13n+k}{10n}\binom{14n+k}{10n}
  \binom{15n}{n+k}\binom{15n}{k}\in\Z.
\end{equation}
Hence $E_n\F_n(1)=\sum_kB_{n,k}$ is an integer, and the same $E_n$
clears every integer polynomial weight in $k$.


\section{The limit and its exponential rate}\label{limit-interface}
The analytic conclusion to use is
\begin{equation}\label{geometric-error}
 r_n(A,B)\longrightarrow\alpha_{A,B},\qquad
 \limsup_{n\to\infty}|r_n(A,B)-\alpha_{A,B}|^{1/n}
 \leq\vartheta:=\frac{|\rho_4|}{|\rho_1|}<1.
\end{equation}
Here $\rho_1,\rho_2,\rho_4$ are the three roots of the exact restricted
characteristic polynomial in Appendix~\ref{limit-details};
$|\rho_1|=|\rho_2|>1>|\rho_4|>0$.
The geometric factor $\vartheta$ and the logarithmic error exponent
$\beta$ are different quantities. If $\vartheta<e^{-\beta}$, the
root-test definition of limsup gives the eventual bound
$|r_n-\alpha_{A,B}|\leq e^{-\beta n}$.
The exact spectral certificate in this handoff gives
\[
 \vartheta<\frac{102}{10^{51}}<e^{-112.8}.
\]
Thus the proposed analytic module supplies $\beta=112.8$ for the
arithmetic comparison; no denominator estimate is used to obtain it.

The proof of \eqref{geometric-error} uses the Meijer--G state only to
identify the limit. Its initial values give the two pairings
$6\alpha_{A,B}$ and $6$. A positive integral proves that this state
decays. The restricted three-dimensional system has only one decaying
mode, so the state selects that mode. The inverse-transpose walk turns
it into the dominant mode. Finally, an exact rank-three intertwining
identity transfers its ratio to the original forward walk.

\subsection{The precise asymptotic statement needed}
To avoid concealing a hypothesis in the phrase ``by asymptotics'',
we isolate the matrix statement used in that transfer. A modal-basis
theorem for rational difference systems must supply the basis in the
next lemma. This is the general asymptotic theorem to formalize or
import; the subsequent transfer is elementary linear algebra.

\begin{lemma}[Dual mode transfer]\label{dual-mode-transfer}
Let $Q(n)\in\operatorname{GL}_3(\mathbb C)$ and
$W_Q(n)=Q(0)\cdots Q(n-1)$. Suppose there are three independent row
solutions, assembled as $B_n=B_0W_Q(n)$, and nonzero scalars $a_i(n)$
such that, for $n\geq1$,
\[
 Y_n=\operatorname{diag}(a_1(n),a_2(n),a_3(n))^{-1}B_nT_n^{-1}
 \longrightarrow Y\in\operatorname{GL}_3(\mathbb C),\quad
 T_n=\operatorname{diag}(1,n,n^2),
\]
\[
 |a_1(n)|^{1/n}\to r,qquad
 |a_2(n)|^{1/n},|a_3(n)|^{1/n}\to R,qquad 0<r<R.
\]
Write $s_0$ for the first row of $B_0$. Let columns $v_n$ satisfy
$T_n^{-T}v_n\to w$, and assume $(Y^{-T}w)_1\ne0$.
For fixed rows $\ell_1,\ell_2$ with $\ell_2s_0^T\ne0$, their ratio
on $W_Q(n)^{-T}v_n$ is eventually defined and tends to
$\alpha=\ell_1s_0^T/(\ell_2s_0^T)$, with
\[
 \limsup_{n\to\infty}
 \left|\frac{\ell_1W_Q(n)^{-T}v_n}
 {\ell_2W_Q(n)^{-T}v_n}-\alpha\right|^{1/n}\leq r/R.
\]
\end{lemma}
\begin{proof}
Put $b_n=Y_n^{-T}T_n^{-T}v_n\to b=Y^{-T}w$.
The factorization of $B_n$ gives the exact identity
\[
 W_Q(n)^{-T}v_n
 =B_0^T\operatorname{diag}(a_1(n)^{-1},a_2(n)^{-1},a_3(n)^{-1})b_n.
\]
For large $n$, $b_{n,1}\ne0$. Dividing by the common nonzero scalar
$b_{n,1}/a_1(n)$ leaves
\[
 s_0^T+\epsilon_n,\qquad
 \epsilon_n=\sum_{i=2}^3
 \frac{a_1(n)b_{n,i}}{a_i(n)b_{n,1}}(B_0)_{i,*}^T,
 \qquad\limsup\|\epsilon_n\|^{1/n}\leq r/R.
\]
The denominator tends to $\ell_2s_0^T\ne0$.
Subtracting $\alpha$ expresses the ratio error as
$(\ell_1-\alpha\ell_2)\epsilon_n/
 (\ell_2s_0^T+\ell_2\epsilon_n)$ and proves both assertions.
\end{proof}

For this application the balanced matrices
$T_nQ(n)T_{n+1}^{-1}$ have rational entries and a finite invertible
limit $Q_\infty$ with distinct eigenvalues. The required modal-basis
result asserts that the rows of $Y$ can be chosen as left eigenvectors
of $Q_\infty$, with the rates above equal to their eigenvalue moduli.
It must be invoked with its exact hypotheses, not inferred from three
numerical eigenvalues. Once that result is supplied, the exponentially
decaying Meijer--G state is a nonzero multiple of the first row solution:
multiplying its expansion by $Y_n^{-1}$ rules out both growing
coordinates. Rescale that row to equal the actual state.

The transported projection, after a scalar that cancels from the ratio,
satisfies $T_n^{-T}v_n\to(1,10,100)^T$. The exact nonzero determinant
\[
 \det[w,Q_\infty^{-T}w,(Q_\infty^{-T})^2w]\ne0
\]
ensures $(Y^{-T}w)_1\ne0$: in an eigenbasis this determinant is the
product of the three coordinates of $w$, a nonzero change-of-basis
determinant, and a Vandermonde determinant in the distinct eigenvalues.
The required initial denominator pairing is $6$, not an assumed
irrationality statement. All trajectory data, integral computations,
and exact transport identities are given in Appendix~\ref{limit-details}.

\section{The harmonic readouts}
To express the numerator readouts, deform the coefficient product by
one auxiliary parameter. With variable lower argument, use
\[
 \binom{M}{x}=\frac{M!}{\Gamma(x+1)\Gamma(M-x+1)}\qquad(M\in\Z_{\geq0}),
\]
while $\binom{x}{M}$ denotes the usual binomial polynomial.
For every integer $j\geq0$, define
\begin{equation}\label{T}
 \mathcal T_{n,j}(\varepsilon)=
 \prod_{u=10}^{12}\binom{j+un-\varepsilon}{10n}
 \binom{15n}{j-n-\varepsilon}
 \binom{15n}{j-2n-\varepsilon},
\end{equation}
\[
 U_{n,j}=\mathcal T_{n,j}(0),\qquad
 V_{n,j}=\mathcal T'_{n,j}(0).
\]
Thus $U_{n,2n+k}=B_{n,k}$ for $0\leq k\leq14n$.
For $j<n$ or $j>17n$, both of the last two factors in \eqref{T}
vanish at zero, so $U_{n,j}=V_{n,j}=0$. All sums below are consequently
finite. The derivative is taken before discarding zero-valued terms:
a single vanishing factor can have a nonzero derivative.

Parameter derivatives produce harmonic differences by
$\psi(m+1)=H_m-\gamma$ \cite[Eq.~5.4.14]{DLMF}.
The two numerator readouts can be written directly as
\begin{align}
 S_{0,n}&=\sum_{j\geq0}U_{n,j}=E_n\F_n(1),\label{S0}\\
 S_{3,n}&=\sum_{j\geq0}\left(U_{n,j}H_{j+2n}^{(3)}
                      +\tfrac12V_{n,j}H_{j+2n}^{(2)}\right),\label{S3}\\
 S_{4,n}&=\sum_{j\geq0}\left(U_{n,j}H_{j+2n}^{(4)}
                      +\tfrac13V_{n,j}H_{j+2n}^{(3)}\right).\label{S4}
\end{align}
In this normalization the forward readout identities are
\begin{equation}\label{walk-readout}
 r_QW_ne_0=\frac{12}{E_n}S_{0,n},\qquad
 r_{A,B}W_ne_0=\frac{12}{E_n}(AS_{3,n}+BS_{4,n}).
\end{equation}
The proposed contiguity derivation is supplied in
Appendix~\ref{arith-walk-section}; its verification is a separate
obligation from the arithmetic of the finite sums.
Under these readout identities, we identify
the original approximants as
\[
 r_n(A,B)=\frac{AS_{3,n}+BS_{4,n}}{S_{0,n}}.
\]
The arithmetic argument that follows is a statement about the explicit
finite sums \eqref{S0}--\eqref{S4}. It does not identify the derivative
observable $e_0^TW_nc_{A,B}$ with the different readout
$r_{A,B}W_ne_0$.

\section{Integer normalization}
Set
\[
 D_n=d_{19n}=\operatorname{lcm}(1,\ldots,19n).
\]
This covers every contributing harmonic index $j+2n$ and will also
clear the parameter derivatives in \eqref{T}.

\begin{proposition}\label{integral}
For every $n\geq1$ and $j\geq0$,
\[
 U_{n,j}\in\Z,\qquad D_nV_{n,j}\in\Z.
\]
Consequently $Z_{i,n}:=6D_n^4S_{i,n}$ is an integer for $i=0,3,4$.
\end{proposition}
\begin{proof}
The assertion for $U_{n,j}$ follows from its binomial product.
For a fixed lower argument, differentiating Vandermonde's identity gives
\[
 \left.\frac{d}{dx}\binom{x}{M}\right|_{x=N}
 =\sum_{h=1}^{M}\binom{N}{M-h}\frac{(-1)^{h-1}}h;
\]
hence $d_M$ clears the derivative at every integer $N$.
For a fixed upper argument, the standard digamma identity gives
\[
 \left.\frac{d}{d\varepsilon}\binom{M}{m-\varepsilon}
 \right|_{\varepsilon=0}
 =\binom{M}{m}(H_m-H_{M-m})\qquad(0\leq m\leq M).
\]
For $m=-r$ or $m=M+r$, $r\geq1$, the derivative instead has absolute
value $M!(r-1)!/(M+r)!$, by the simple zeros of reciprocal Gamma
\cite[\S5.2]{DLMF}. This equals the absolute value of the derivative of
$\binom{x}{M+r}$ at $x=r-1$, so it is cleared by $d_{M+r}$.

For $0\leq j\leq17n$, the polynomial factors in \eqref{T} have degree
$10n$. The fixed-upper factors have upper argument $15n$, and their
lower arguments lie at most $2n$ outside the usual range. Thus $D_n$
clears every differentiated factor. The product rule proves
$D_nV_{n,j}\in\Z$; outside this range the whole derivative vanishes.
Finally, \eqref{harmonic-integrality} applied to \eqref{S3}--\eqref{S4}
proves $6D_n^4S_{i,n}\in\Z$. One power of $D_n$ clears $V_{n,j}$,
and the remaining three clear the highest accompanying harmonic order.
The factor $6$ clears the fixed multipliers $1/2$ and $1/3$.
\end{proof}

\section{Factorial permutations and a common divisor}
Let $b=(b_0,\ldots,b_4)$ run over the permutations of $(-2,-1,0,1,2)$.
Write
\[
 L_b=(10+b_0,11+b_1,12+b_2),\qquad
 J_b=(16-b_3,17-b_4),
\]
and define the alternative coefficient product
\[
 \mathcal T^b_{n,j}(\varepsilon)=
 \prod_{i=0}^2\binom{j+(10+i)n-\varepsilon}{(10+i+b_i)n}
 \binom{(16-b_3)n}{j-b_3n-\varepsilon}
 \binom{(17-b_4)n}{j-b_4n-\varepsilon}.
\]
Expansion into Gamma factors shows that
\begin{equation}\label{permutation}
 \mathcal T^b_{n,j}(\varepsilon)=\lambda_{n,b}\mathcal T_{n,j}(\varepsilon),
 \qquad
 \lambda_{n,b}=
 \frac{(10n)!^3\prod_{e\in J_b}(en)!}
 {(15n)!^2\prod_{\ell\in L_b}(\ell n)!}.
\end{equation}
Indeed the three varying numerator factors are unchanged, and the five
factors $\Gamma(j-b_i n+1-\varepsilon)$ in the denominator are merely
permuted. The two other varying factors are
$\Gamma(16n-j+1+\varepsilon)$ and
$\Gamma(17n-j+1+\varepsilon)$.
The scalar $\lambda_{n,b}$ is exactly the ratio of the constant
factorials in the two products.

The proof of Proposition~\ref{integral} applies to every permutation.
The three polynomial degrees are at most $14n$. For each fixed-upper
factor, writing its arguments as $M$ and $m$, all of $M,m,M-m$ are
at most $19n$ when $0\leq j\leq17n$; both the interior derivative and
the exterior derivative are therefore cleared by $D_n$.
For larger $j$, the value and derivative vanish by
\eqref{permutation}. It follows that
\[
 \lambda_{n,b}U_{n,j}\in\Z,\qquad
 D_n\lambda_{n,b}V_{n,j}\in\Z,
 \qquad \lambda_{n,b}Z_{i,n}\in\Z\quad(i=0,3,4).
\]
Thus each permutation yields an integer multiple of the same rational
readout vector, and the multiplier is known explicitly.

For a prime $p$ let $v_p$ denote its valuation, with $v_p(0)=+\infty$.
Define
\begin{equation}\label{divisor}
 e_p(n)=\max\left(0,\max_b\{-v_p(\lambda_{n,b})\}\right),\qquad
 \mathcal G_n=\prod_{p\leq19n}p^{e_p(n)}.
\end{equation}
Legendre's formula makes the exponents explicit:
\[
 -v_p(\lambda_{n,b})=\sum_{h\geq1}\left[
 2\left\lfloor\frac{15n}{p^h}\right\rfloor
 -3\left\lfloor\frac{10n}{p^h}\right\rfloor
 -\sum_{e\in J_b}\left\lfloor\frac{en}{p^h}\right\rfloor
 +\sum_{\ell\in L_b}\left\lfloor\frac{\ell n}{p^h}\right\rfloor\right].
\]
Only finitely many terms are nonzero. The maximum over permutations
is taken after the sum over prime powers.

\begin{proposition}\label{gcd}
For every $n\geq1$,
\[
 \mathcal G_n\mid\gcd(Z_{0,n},Z_{3,n},Z_{4,n}).
\]
Consequently, for every integer pair $A,B$,
\[
 \mathcal G_n\mid\gcd(Z_{0,n},AZ_{3,n}+BZ_{4,n}).
\]
If $r_n(A,B)=p_n/q_n$ is reduced with $q_n>0$, then
\begin{equation}\label{height}
 q_n\leq\frac{6D_n^4E_n\F_n(1)}{\mathcal G_n}.
\end{equation}
\end{proposition}
\begin{proof}
For each permutation, integrality of $\lambda_{n,b}Z_{i,n}$ gives
$v_p(Z_{i,n})\geq-v_p(\lambda_{n,b})$. Since $Z_{i,n}$ is itself an
integer, $v_p(Z_{i,n})\geq e_p(n)$. This proves the first assertion;
the second follows by integer linear combination.
Moreover, $Z_{0,n}=6D_n^4\sum_kB_{n,k}>0$.
Dividing numerator and denominator by $\mathcal G_n$ and then reducing
the resulting fraction proves \eqref{height}.
\end{proof}

The same reasoning applied coefficientwise shows
$\mathcal G_n\mid B_{n,k}$ for every $k$. Since $B_{n,0}=E_n$,
$E_n/\mathcal G_n$ is an integer and clears every polynomially weighted
sum in \eqref{integer-weight}. The harmonic readouts require the
additional lcm normalization established in Proposition~\ref{integral}.

\section{The error-versus-height criterion}
Let $\alpha_{A,B}=A\zeta(3)+B\zeta(4)$ and set
\[
 u_n=Z_{0,n}/\mathcal G_n,\qquad
 v_n(A,B)=(AZ_{3,n}+BZ_{4,n})/\mathcal G_n.
\]
Proposition~\ref{gcd} gives integers $u_n>0$, $v_n(A,B)$ with
$r_n(A,B)=v_n(A,B)/u_n$.

\begin{corollary}
Fix $(A,B)\in\Z^2\setminus\{(0,0)\}$. Suppose there is an infinite
set $\mathcal N$ of depths on which
\begin{equation}\label{assumed-error}
 0<|\alpha_{A,B}-r_n(A,B)|,\qquad
 u_n|\alpha_{A,B}-r_n(A,B)|\longrightarrow0.
\end{equation}
Then $\alpha_{A,B}$ is irrational. If this hypothesis holds for every
fixed nonzero integer pair, then $1,\zeta(3),\zeta(4)$ are linearly
independent over $\mathbb Q$.
\end{corollary}
\begin{proof}
If $\alpha_{A,B}=a/b$ with integers $a,b$ and $b>0$, then for every
$n\in\mathcal N$,
\[
 u_n|\alpha_{A,B}-r_n(A,B)|
 =\frac{|au_n-bv_n(A,B)|}{b}\geq\frac1b,
\]
since the numerator is a nonzero integer. This contradicts
\eqref{assumed-error}. A nontrivial rational relation among
$1,\zeta(3),\zeta(4)$ would make one such nonzero integer combination
rational, proving the second assertion.
\end{proof}

Quantitatively, if the assumed approximation bound and the arithmetic
height bound on the same depths are
\[
 |\alpha_{A,B}-r_n(A,B)|\leq e^{-\beta n+o(n)},\qquad
 u_n\leq e^{\sigma n+o(n)},\qquad \beta>\sigma>0,
\]
with nonzero error, then \eqref{assumed-error} holds. For the reduced
denominators, $q_n\leq u_n$ and
\[
 \liminf_{\substack{n\to\infty\\n\in\mathcal N}}
 \left(-1-\frac{\log|\alpha_{A,B}-p_n/q_n|}{\log q_n}\right)
 \geq\frac{\beta}{\sigma}-1>0.
\]
The denominators tend to infinity on these depths because the limit
is irrational. The analytic estimate is the output of Section~\ref{limit-interface}, subject to its stated modal-basis dependency. In particular,
accepting an approximation exponent such as $112.8$ still requires a
height exponent strictly below it for the divisor actually established
in \eqref{divisor}. No such numerical comparison is claimed by the
permutation divisibility argument alone.


\section{The stronger divisor and the quantitative conclusion}
The coefficientwise permutation divisor $\mathcal G_n$ above should not
be confused with the stronger divisor $G_n$ in the appendix. The latter
uses eight factorial pairings, the upper harmonic ranges, and one
cancellation over all residues modulo a prime. Its complete arithmetic
argument is reproduced in the appendix, including the finite floor
partition and the rational logarithm bounds.

To translate its notation, its rational function is denoted $F_n(t)$,
whereas our terminating hypergeometric polynomial is $\F_n(z)$.
The appendix uses
\[
 C_n=\frac{(15n)!^2}{(10n)!^3},\qquad
 \rho_n=(-1)^n\frac{(12n)!(13n)!}
 {n!(2n)!(3n)!(4n)!(15n)!}.
\]
They satisfy $E_n=C_n|\rho_n|$. Direct coefficient comparison gives
\[
 U_{n,j}=(-1)^n C_na_j,\quad V_{n,j}=(-1)^n C_nb_j,
 \qquad (S_{0,n},S_{3,n},S_{4,n})
 =(-1)^n C_n(Q_n,P_{3,n},P_{4,n}).
\]
Thus the appendix's integer vector $X_n$ equals
$(-1)^n(Z_{0,n},Z_{3,n},Z_{4,n})$: both normalizations have exactly the
same common divisors. In particular $G_n$ divides all three $Z_{i,n}$.
Set
\[
 \widetilde u_n=Z_{0,n}/G_n>0,\qquad
 \widetilde v_n(A,B)=(AZ_{3,n}+BZ_{4,n})/G_n.
\]
Both are integers. The height argument in the appendix bounds the
\emph{unreduced} denominator $\widetilde u_n$ itself:
\[
 \limsup_{n\to\infty}\frac{\log\widetilde u_n}{n}
 \leq65.46+76-(S_8+9)<112.5,
\]
where $S_8>20.004$ is its explicit rational sum. Consequently the
reduced denominator is at most $\widetilde u_n$, and
$\widetilde u_n\leq e^{112.5n}$ for all sufficiently large $n$.

\begin{theorem}[Conditional application]\label{conditional-completion}
Assume the finite-sum readout identification and the recurrence
identification described below, and the analytic module of
Section~\ref{limit-interface}. Equivalently, its rate conclusion may
be taken as the following explicit input. For each fixed nonzero
$(A,B)\in\Z^2$, assume
\begin{equation}\label{rate-input}
 |A\zeta(3)+B\zeta(4)-r_n(A,B)|
 \leq\exp(-\beta n+o(n)),\qquad \beta>112.5.
\end{equation}
The $o(n)$ term may depend on $A,B$. With the stronger arithmetic bound
of the appendix, every such combination is irrational, and hence
$1,\zeta(3),\zeta(4)$ are linearly independent over $\Q$.
If $\beta=112.8=564/5$, their reduced approximants satisfy
\[
 \liminf_{n\to\infty}\delta_n(A,B)\geq\frac{112.8}{112.5}-1
 =\frac1{375}>0.
\]
After irrationality is established, every rational approximation has
nonzero error and its reduced denominator tends to infinity.
\end{theorem}
\begin{proof}
The nonvanishing argument below supplies infinitely many depths with
nonzero error for each fixed pair. On these depths the two bounds imply
\[
 0<\widetilde u_n|A\zeta(3)+B\zeta(4)-r_n(A,B)|
 \leq e^{-(\beta-112.5)n+o(n)}\longrightarrow0.
\]
If the limit were $a/b$, with $a,b\in\Z$ and $b>0$, the left-hand
side would be at least $1/b$, since
$a\widetilde u_n-b\widetilde v_n$ is a nonzero integer.
This proves irrationality. Clearing denominators in any rational
relation among $1,\zeta(3),\zeta(4)$ proves the independence assertion.
Finally $q_n\leq\widetilde u_n$ gives the stated delta bound.
\end{proof}
Only the exact hypothesis $\beta>112.5$ is used here. A bound with
$\beta=112$ alone is insufficient for this argument. The numerical
margin $1/375$ is a conditional lower bound, not a measured delta.

\subsection{Nonvanishing for every fixed integer pair}
The auxiliary forward Meijer--G matrix $M_G(n)$ in this recurrence
argument has the parameters in Appendix~\ref{arith-walk-section}.
It differs from the dual Meijer--G matrix $G(n)$ used in the limit
argument. Their initial matrices and scalar normalizations must be
kept distinct.
Here is the precise recurrence argument needed, including the case
where the target constant might be zero. Let $M_G(n)$ be the native
Meijer--G matrix specified in Appendix~\ref{arith-walk-section}.
The certificate file gives a three-dimensional companion $C^{(G)}(n)$
and a $4\times3$ rational matrix $U(n)$ with
\[
 M_G(n)U(n+1)=U(n)C^{(G)}(n),\quad U(n)e_0=e_0,\qquad
 C^{(G)}(n)=\begin{pmatrix}0&0&c_0(n)\\1&0&c_1(n)\\0&1&c_2(n)\end{pmatrix}.
\]
Its rational-function identities, absence of poles at $n\geq0$, and
$c_0(n)\ne0$ at those indices are certificate obligations. The verifier
checks the numerator of $c_0$ and the three denominators by factoring:
each factor has a nonzero constant coefficient and coefficients of one
sign. The full coefficient data are included, not just sample values.
The readout identification needed here is
\[
 (Q_n,P_{3,n},P_{4,n})^T=(-1)^n R W_n^{(G)}e_0,\qquad
 R=\begin{pmatrix}1&0&0&0\\0&0&-1/2&1/4\\0&0&0&-1/6\end{pmatrix}.
\]
It is the same bridge as in the appendix, in the equivalent Meijer--G
coordinates. It must be proved for all $n$, including the initial
values $(Q_0,P_{3,0},P_{4,0})=(1,0,0)$.

Put $S=RU(0)$. The saved exact certificate gives $\det S\ne0$;
one short check is
\[
 S\bmod101=\begin{pmatrix}1&1&56\\0&56&10\\0&76&19\end{pmatrix},
 \qquad\det S\equiv1\pmod{101}.
\]
All denominators of its entries are prime to $101$.
For any real $\alpha$ and integers $A,B$, let
\[
 z_n=(-1)^n(\alpha Q_n-AP_{3,n}-BP_{4,n}).
\]
Telescoping the coboundary identity shows that $z_n$ is the first
coordinate of a row companion solution with initial row
$(\alpha,-A,-B)S$. Its first three scalar values are exactly that row,
and it obeys
\[
 z_{n+3}=c_0(n)z_n+c_1(n)z_{n+1}+c_2(n)z_{n+2}.
\]
If $z_n$ were eventually zero, three consecutive zeros and $c_0(n)\ne0$
would force the preceding value to vanish. Backward induction would
force $z_0=z_1=z_2=0$. Invertibility of $S$ would then give
$\alpha=A=B=0$. This is impossible when $(A,B)\ne(0,0)$.
Taking $\alpha=A\zeta(3)+B\zeta(4)$ proves the required infinitely many
nonzero errors, without assuming that $\alpha$ is nonzero.

\paragraph{Order of verification.}
The arithmetic theorem concerns the explicit finite sums. The readout
identity connects those sums with the forward walk. The recurrence
excludes an eventually exact approximation. The limit argument identifies
the zeta combination and supplies its error exponent. These statements
must be connected for the same sequences before applying
Proposition~\ref{conceptual-criterion}. The exact matrix checks are
included as finite certificates; the general modal-basis theorem remains
a separate analytic dependency.

\clearpage\appendix
\section{The limit construction in detail}\label{limit-details}

We write $(\mathcal M,x_0,v)_F$ for the forward scheme and
$(\mathcal M,x_0,v)_D$ for the dual scheme of a CMF $\mathcal M$ with
starting point $x_0$ and direction $v$.

\paragraph{Analytic dependency.}
The limit theorem below uses the modal-basis assertion specified in
Section~\ref{limit-interface}. Its integral estimates and exact matrix
identities are supplied here. The proof of that general difference-system
assertion is a separate formalization task; the characteristic roots
alone are not a substitute for it.

\begin{theorem}
\label{thm:z34-limit}
Let $A,B\in\mathbb Z$, $x_0=(0,0,1,1,1;1,1,1,1)$ and
$v_F=(-14,-15,13,14,12;1,2,3,4)$.
The forward scheme
\[
\bigl({}_5F_4(1),\,x_0,\,v_F\bigr)_F
\]
has initial condition and projection vector
\[
I_F=\begin{pmatrix}0&0&-6A&3A-2B\\12&0&0&0\end{pmatrix},
\qquad v=\begin{pmatrix}1\\0\\0\\0\end{pmatrix}.
\]
Writing $W_F(N)=F_{x_0,v_F}(0)\cdots F_{x_0,v_F}(N-1)$, the ratio
\[
R_N=\frac{(I_FW_F(N)v)_1}{(I_FW_F(N)v)_2}
\]
converges to $A\zeta(3)+B\zeta(4)$, with exponential error bound
\[
\limsup_{N\to\infty}|R_N-A\zeta(3)-B\zeta(4)|^{1/N}
\leq\vartheta:=\frac{|\rho_4|}{|\rho_1|}
\approx1.019\cdot10^{-49},
\]
where the roots are those of Proposition~\ref{prop:z34-rates}.
\end{theorem}
Unlike the preceding case, we cannot identify the full dual Meijer G
scheme with this forward scheme by an invertible rational coboundary.
Instead, a rank-three intertwining matrix connects the selected walks.
We show that this connection preserves the dominant dual contribution
and the initial-state identity determining the limit.

\begin{proposition}
\label{prop:z34-coboundary}
Let $y_0=(0,0,1,1,2;0,0,0,0,0)$ and
$v_G=(-16,-17,11,12,10;-2,-1,0,1,2)$.
The dual scheme $(G^{5,2}_{5,5}(1),y_0,v_G)_D$, restricted to its first
three coordinates, is coboundary equivalent, up to scalar normalization,
to the forward scheme $({}_5F_4(1),x_0,v_F)_F$ on a three-dimensional
invariant space $\operatorname{im}E(n)$.
\end{proposition}

\begin{proof}

The matrices and scalar factors are explicit:
\[
 E(n)=\begin{pmatrix}
 1&0&0\\
 0&-n(450n+29)&-n^2(101250n^2+12975n+406)\\
 0&1&0\\0&0&1
 \end{pmatrix},\quad
 L=\begin{pmatrix}1&0&0&0\\0&0&1&0\\0&0&0&1\end{pmatrix}.
\]
Write $H(n)=(h_{ij})_{0\leq i,j\leq2}$, where
\begin{align*}
 h_{00}&=-6n^2(10n+1)(12n+1)(193n+13),&
 h_{01}&=-n(12n+1)(2550n^2+523n+23),\\
 h_{02}&=(12n+1)(450n^2+41n+1),&
 h_{10}&=(10n+1)(450n^2+21n-1),\\
 h_{11}&=(-25n-2)(18n-1),&h_{12}&=-1,\\
 h_{20}&=-20n-2,&h_{21}&=2,\qquad h_{22}=0.
\end{align*}
The positive scalar $R(n)=\mathfrak a(n)/\mathfrak b(n)$ is given by
\begin{align*}
 \mathfrak a(n)&=1125(n+1)^3(2n+1)(3n+1)(3n+2)
 \prod_{j=1}^4(5n+j)
 \prod_{j\in\{1,2,4,7,8,11,13,14\}}(15n+j),\\
 \mathfrak b(n)&=416(6n+1)(6n+5)
 \prod_{j\in\{1,5,7,11\}}(12n+j)
 \prod_{j=1}^{12}(13n+j).
\end{align*}
In particular, $\det H(n)=2\prod_{j=8}^{12}(jn+1)\ne0$ for $n\geq0$.
The attached symbolic certificate verifies the following identities
as rational functions, not only at finitely many values:
\[
F_{x_0,v_F}(n)E(n+1)=E(n)\widehat F(n),
\qquad
\widehat F(n)=\frac{\mathfrak a(n)}{\mathfrak b(n)}H(n)Q(n)^{-T}H(n+1)^{-1}.
\]
Here $Q(n)$ is the leading $3\times3$ block of $G_{y_0,v_G}(n)$,
$\widehat F(n)=LF_{x_0,v_F}(n)E(n+1)$, and $L$ selects the first, third
and fourth coordinates, with $LE(n)=I_3$.
The restriction of the dual Meijer G matrix follows from
\[
G_{y_0,v_G}(n)=\begin{pmatrix}Q(n)&b(n)\\0&-1\end{pmatrix},
\qquad
G_{y_0,v_G}(n)^{-T}
=\begin{pmatrix}Q(n)^{-T}&0\\b(n)^TQ(n)^{-T}&-1\end{pmatrix}.
\]
Thus a row $(u,0)$ is mapped to $(uQ(n)^{-T},0)$, so the first
three coordinates form an invariant space for the dual walk.
Put
\[
J=\begin{pmatrix}1&0&0&0\\0&1&0&0\\0&0&1&0\end{pmatrix},
\qquad K(n)=E(n)H(n)J,
\qquad R(n)=\frac{\mathfrak a(n)}{\mathfrak b(n)}.
\]
The matrix $K(n)$ has rank three, and the preceding identities imply
\[
F_{x_0,v_F}(n)K(n+1)=R(n)K(n)G_{y_0,v_G}(n)^{-T}.
\]
Consequently, writing $W_G(N)=G_{y_0,v_G}(0)\cdots
G_{y_0,v_G}(N-1)$, we obtain the exact walk identity
\[
W_F(N)K(N)
=\left(\prod_{n=0}^{N-1}R(n)\right)K(0)W_G(N)^{-T}.
\]
This identity does not require $K(n)$ to be invertible.
\end{proof}
Along the trajectory define
\[
\mathcal{G}_n(z)=G_{5,5}^{5,2}\!\left(
 z\;\middle|\;
 \begin{matrix}
 -16n,-17n;\ 1+11n,1+12n,2+10n\\
 -2n,-n,0,n,2n
 \end{matrix}\right),
\qquad \theta=z\,\frac{d}{dz}.
\]

\begin{proposition}
\label{prop:z34-initial-values}
Let
\[
\mathcal{S}_G(n)=\bigl(
\mathcal{G}_n(1),\ (\theta\mathcal{G}_n)(1),\
(\theta^2\mathcal{G}_n)(1),\ (\theta^3\mathcal{G}_n)(1)\bigr),
\]
\[
I_G=\begin{pmatrix}2B&-3A-2B&3A&0\\0&0&6&0\end{pmatrix}.
\]
Then
\[
I_G\mathcal{S}_G(0)^T
=6\begin{pmatrix}A\zeta(3)+B\zeta(4)\\1\end{pmatrix}.
\]
\end{proposition}

\begin{proof}
The Mellin--Barnes integral in the proof of
Proposition~\ref{prop:z34-growth}, with $n=0$, and the reflection formula
\cite[Eq.~5.5.3]{DLMF} give
\[
\mathcal{G}_0(z)=\frac{1}{2\pi i}\int_{\frac12+i\mathbb R}
\frac{\pi^2z^{-s}}{\sin^2(\pi s)s^3(s+1)}\,ds.
\]
Applying $\theta^k$ multiplies the integrand by $(-s)^k$.
At $z=1$, close the contour to the right along lines
$\operatorname{Re}s=m+\tfrac12$. For $0\leq k\leq3$ their integrals are
$O(m^{k-4})\to0$; the horizontal integrals vanish exponentially.
The poles at positive integers are double, and the clockwise orientation gives
\[
(\theta^k\mathcal{G}_0)(1)
=(-1)^{k+1}\sum_{\ell=1}^{\infty}
\left[\frac{d}{ds}\frac{1}{s^{3-k}(s+1)}\right]_{s=\ell}.
\]
Partial fractions give
\[
\mathcal{S}_G(0)=\bigl(
1-2\zeta(3)+3\zeta(4),\ 1-2\zeta(3),\ 1,\ 1-\zeta(2)\bigr),
\]
and multiplication by $I_G$ gives the claimed identity.
\end{proof}

\begin{proposition}
\label{prop:z34-rates}
The trajectory matrix $G_{y_0,v_G}(n)$ has exponential growth rates given by
two conjugate dominant roots
\[
\rho_{1,2}\approx(2.803\pm6.487i)\cdot10^{25},
\qquad \rho_3=-1,
\qquad \rho_4\approx7.203\cdot10^{-24}.
\]
No factorial normalization is needed. The characteristic roots at infinity
of the restricted matrix $Q(n)$ are $\rho_1,\rho_2,\rho_4$.
\end{proposition}

\begin{proof}
The roots are those of the characteristic polynomial at infinity, extracted
in the supplied code. The factor corresponding to $\rho_3=-1$ is removed
by the block restriction in Proposition~\ref{prop:z34-coboundary}.
\end{proof}


\paragraph{Exact spectral data and a rational rate certificate.}
The restricted characteristic polynomial is
\[
 \chi(X)=X^3-aX^2+bX-c,
\]
where
\begingroup\scriptsize
\begin{gather*}
 a=\frac{83220721427294106907399200607730930785765129529801897078934559}{1484577825190137274336280508756918272}\,\\
 b=\frac{731334801243107466089356596505717853275074571608716447379201100334794667003154611901941573810757}{146466868948884247884332915718530654652342544}\,\\
 c=\frac{76623012348334313797938377862237790794460627983667564193521662995163373351121226574589853696}{2130520759189220333584634002331004844478493702212809977860478179}\,
\end{gather*}
\endgroup
Its discriminant is negative. Put $u=7203/10^{27}$ and
$\varepsilon=102/10^{51}$. Exact rational arithmetic gives
\[
 c>0,\qquad\chi(0)<0<\chi(u),\qquad u^3<c\varepsilon^2.
\]
There is therefore exactly one real root $r$ and $0<r<u$; the other
two roots are a nonreal conjugate pair. Their common modulus is
$\sqrt{c/r}$ by the product of the three roots. Consequently
\[
 \vartheta^2=\frac{r^3}{c}<\varepsilon^2,
 \qquad \vartheta<\varepsilon<e^{-112.8}.
\]
The last inequality is certified by dyadic reduction and the rational
series bounds for $\log$ given later in the height proof.
The exact limiting matrix, characteristic coefficients, nonzero cyclic
determinant, and rational log bound are recorded in
\texttt{certificates/spectral\_data.json}. The root decimals are only
illustrations; the inequalities above are the certificate to formalize.

\begin{proposition}
\label{prop:z34-growth}
Let
\[
\eta=\frac{16^{16}15^{30}2^2 3^3 4^4}{13^{13}12^{12}30^{30}}
\left(\frac32\right)^{55}\approx8.517\cdot10^{-4}<1.
\]
Then
\[
(\theta^k\mathcal{G}_n)(1)=O(n^k\eta^n),
\qquad k=0,1,2,3.
\]
\end{proposition}

\begin{proof}
Put $L=14n+1$ and $\mathfrak c_n=(14n)!(16n)!/[(11n)!^2(8n+1)!]$.
The Mellin--Barnes representation \cite[Eq.~16.17.1]{DLMF} is
\[
\mathcal{G}_n(z)=\frac{1}{2\pi i}
\int_{2n+\frac12-i\infty}^{2n+\frac12+i\infty}
\frac{
 \Gamma(16n+1-s)\Gamma(17n+1-s)
}{
 \Gamma(s+11n+1)\Gamma(s+12n+1)\Gamma(s+10n+2)
}\prod_{j=-2}^{2}\Gamma(s+jn)\,z^{-s}\,ds.
\]
The contour is traversed from bottom to top. Its left pole families
lie in $s\leq2n$ and its right pole families in $s\geq16n+1$,
so it separates them for every $n$.
With $s=2n+w$, the contour becomes the fixed line
$\mathcal C=\tfrac12+i\mathbb R$.
This is a change of variable; no poles are crossed.
The integral converges absolutely for $z>0$.

Apply Euler's beta integral \cite[Eq.~5.12.1]{DLMF} to
\[
\frac{\Gamma(w+2n)}{\Gamma(w+13n+1)},\qquad
\frac{\Gamma(w+3n)}{\Gamma(w+14n+1)},\qquad
\frac{\Gamma(w+4n)}{\Gamma(w+12n+2)}.
\]
This gives
\[
\begin{aligned}
\mathcal{G}_n(z)
&=\frac{z^{-2n}}{2\pi i(11n)!^2(8n+1)!}
\int_{\mathcal C}
\Gamma(w)\Gamma(w+n)\Gamma(14n+1-w)\Gamma(15n+1-w)z^{-w}\\
&\quad\times\left[\int_{[0,1]^3}
x^{w+2n-1}y^{w+3n-1}t^{w+4n-1}
(1-x)^{11n}(1-y)^{11n}(1-t)^{8n+1}
\,dx\,dy\,dt\right]dw.
\end{aligned}
\]
All beta parameters have positive real parts on $\mathcal C$.
Apply Euler's integral once more:
\[
\Gamma(w+n)\Gamma(15n+1-w)
=(16n)!\int_0^1 r^{w+n-1}(1-r)^{15n-w}\,dr.
\]
Absolute convergence permits interchanging the integrals.
Writing $u=xyt$, the remaining integral on the same contour is
\cite[Eq.~5.13.1]{DLMF}
\[
\frac{1}{2\pi i}\int_{\mathcal C}
\Gamma(w)\Gamma(14n+1-w)
\left(\frac{z(1-r)}{ur}\right)^{-w}dw
=(14n)!\left(\frac{ur}{z(1-r)+ur}\right)^{14n+1}.
\]
Its argument is positive and its contour condition
$0<\tfrac12<14n+1$ holds for every $n$.
Consequently, we obtain the positive integral
\[
\mathcal{G}_n(z)=z^{-2n}\mathfrak c_n\int_{[0,1]^4}
\frac{
x^{16n}y^{17n}t^{18n}r^{15n}
(1-x)^{11n}(1-y)^{11n}(1-t)^{8n+1}(1-r)^{15n}
}{[z(1-r)+xytr]^{14n+1}}
\,dx\,dy\,dt\,dr.
\]
In particular, $z=1$ can be substituted directly.
Since $1-r+xytr\geq xyt$, for $n\geq1$ Euler's integral gives
\[
\begin{aligned}
0<\mathcal{G}_n(1)
&\leq \mathfrak c_n\int_{[0,1]^4}
x^{2n-1}y^{3n-1}t^{4n-1}r^{15n}
(1-x)^{11n}(1-y)^{11n}(1-t)^{8n+1}(1-r)^{15n}
\,dx\,dy\,dt\,dr\\
&=\frac{(16n)!(15n)!^2(2n-1)!(3n-1)!(4n-1)!}
{(13n)!(12n+1)!(30n+1)!}
=:\mathfrak d_n.
\end{aligned}
\]
The factorial ratios satisfy, for $n\geq2$,
\[
\frac{D_{n+1}}{\mathfrak d_n}
\leq
\frac{16^{16}15^{30}2^2 3^3 4^4}{13^{13}12^{12}30^{30}}
\left(\frac{n+1}{n}\right)^{55}
\leq
\frac{16^{16}15^{30}2^2 3^3 4^4}{13^{13}12^{12}30^{30}}
\left(\frac32\right)^{55}
=:\eta\approx8.517\cdot10^{-4}<1.
\]
Here the $55$ numerator factors are bounded above using $n+1$,
and the $55$ denominator factors are bounded below using $n$.
Iterating gives, for a constant $C>0$ independent of $n$,
\[
0<\mathcal{G}_n(1)\leq \mathfrak d_n\leq C\eta^n\longrightarrow0.
\]

Finally, $\theta^k$ inserts the factor $(-s)^k$ in the Mellin--Barnes
integrand. For $k\leq3$, these integrals
remain absolutely convergent at $z=1$ on the same contour.
Put $a=2n$. After $s=a+w$, the factor is
$(-1)^k(w+a)^k$.
For $n\geq1$ and $k\leq3$, the polynomial $(w+a)^k$ has an expansion
in the rising factorials $(w)_j$ with nonnegative coefficients.
Each $(w)_j$ replaces $\Gamma(w)$ by $\Gamma(w+j)$; the same Mellin
inversion then multiplies the positive integrand at $z=1$ by
$(L)_j[(1-r)/(1-r+ur)]^j\leq(L)_j$.
For clarity, the expansions for the only required orders are
\begin{align*}
 (w+a)^0&=1,\\
 (w+a)^1&=(w)_1+a,\\
 (w+a)^2&=(w)_2+(2a-1)(w)_1+a^2,\\
 (w+a)^3&=(w)_3+(3a-3)(w)_2+(3a^2-3a+1)(w)_1+a^3.
\end{align*}
All coefficients are nonnegative when $a=2n\geq2$.
Evaluating the polynomial expansion at $w=L$ bounds the combined
factor by $(a+L)^k$. Hence
\[
|(\theta^k\mathcal{G}_n)(1)|
\leq
(16n+1)^k\mathfrak d_n
\leq C(16n+1)^k\eta^n\longrightarrow0,
\qquad 0\leq k\leq3.
\]
Since $16n+1=O(n)$, the claim follows.
\end{proof}

\begin{proof}[Proof of Theorem~\ref{thm:z34-limit}]
Assume the modal-basis theorem with the hypotheses stated in
Section~\ref{limit-interface}. By Proposition~\ref{prop:z34-growth}, every component of the
state vector of $(G^{5,2}_{5,5}(1),y_0,v_G)_F$ decays exponentially, so by
Proposition~\ref{prop:z34-rates} it does not grow at any of the rates
$|\rho_1|,|\rho_2|,|\rho_3|$; hence the state is the recessive solution.
Write $\widehat{\mathcal S}_G(n)$ for its first three components and
$\widehat I_G$ for the first three columns of $I_G$.
The state relation and the block form in
Proposition~\ref{prop:z34-coboundary} give
\[
\mathcal S_G(n+1)=\mathcal S_G(n)G_{y_0,v_G}(n),
\qquad
\widehat{\mathcal S}_G(n+1)=\widehat{\mathcal S}_G(n)Q(n).
\]
Hence $\widehat{\mathcal S}_G(n)$ is the recessive solution of $Q(n)$
obtained from the actual Meijer G state. Since $I_G=(\widehat I_G\;0)$,
\[
\widehat I_G\widehat{\mathcal S}_G(0)^T
=I_G\mathcal S_G(0)^T
=6\begin{pmatrix}A\zeta(3)+B\zeta(4)\\1\end{pmatrix}.
\]
The block form also gives the exact restriction
\[
I_GW_G(N)^{-T}
=\bigl(\widehat I_GW_Q(N)^{-T}\;0\bigr),
\qquad W_Q(N)=Q(0)\cdots Q(N-1).
\]
Thus the initial conditions already exclude the $-1$ mode; the
rank-three connection transfers this restriction to the forward walk.
Using the modal-basis statement and Lemma~\ref{dual-mode-transfer},
with the projection condition verified below, the dual scheme restricted to
$Q(n)$, with initial condition $\widehat I_G$, converges to
$(I_G\mathcal{S}_G(0)^T)_1/(I_G\mathcal{S}_G(0)^T)_2$.
Proposition~\ref{prop:z34-initial-values} proves that this ratio is
$A\zeta(3)+B\zeta(4)$.

At the source point, the coboundary of
Proposition~\ref{prop:z34-coboundary} satisfies
\[
E(0)=\begin{pmatrix}1&0&0\\0&0&0\\0&1&0\\0&0&1\end{pmatrix},
\qquad H(0)=\begin{pmatrix}0&0&1\\-1&2&-1\\-2&2&0\end{pmatrix},
\]
\[
I_FE(0)H(0)=2\widehat I_G.
\]
Equivalently, $I_FK(0)=2I_G$.
Consequently,
\[
I_FE(0)H(0)\widehat{\mathcal S}_G(0)^T
=12\begin{pmatrix}A\zeta(3)+B\zeta(4)\\1\end{pmatrix}.
\]
Put $e_0=(1,0,0)^T$ and $\Delta(n)=\prod_{j=8}^{12}(jn+1)$.
The supplied code verifies
\[
H(n)^{-1}e_0=\frac{1}{\Delta(n)}
\begin{pmatrix}1\\10n+1\\(10n+1)^2\end{pmatrix},
\qquad K(n)J^TH(n)^{-1}e_0=v.
\]
Thus the exact walk identity gives
\[
I_FW_F(N)v
=2\left(\prod_{n=0}^{N-1}\frac{\mathfrak a(n)}{\mathfrak b(n)}\right)
\widehat I_G W_Q(N)^{-T}H(N)^{-1}e_0.
\]
The common scalar cancels in the ratio. For
$T_n=\operatorname{diag}(1,n,n^2)$ and
$Q_\infty=\lim_{n\to\infty}T_nQ(n)T_{n+1}^{-1}$, the transported
projection satisfies $\Delta(n)T_n^{-1}H(n)^{-1}e_0\to w=(1,10,100)^T$.
The supplied code verifies exactly that $Q_\infty$ has distinct
eigenvalues and
\[
\det\bigl[w,\ Q_\infty^{-T}w,\ (Q_\infty^{-T})^2w\bigr]\ne0.
\]
Thus $w$ has a nonzero coefficient along every eigenvector in an
eigenbasis of $Q_\infty^{-T}$, including the eigenvector for the unique
dominant dual root $\rho_4^{-1}$. The modal-basis statement in Section~\ref{limit-interface} and
Lemma~\ref{dual-mode-transfer} now give the asserted contribution and
rate. The balancing matrices and scalar $\Delta(n)$ introduce only polynomial
factors and a scalar common to numerator and denominator.
The denominator's pairing with the initial state is nonzero, so
this gives $R_N\to A\zeta(3)+B\zeta(4)$.
Finally, Proposition~\ref{prop:z34-rates} gives the exponential error
factor $\vartheta=|\rho_4|/|\rho_1|$, since $|\rho_1|=|\rho_2|$.
\end{proof}



\section*{Detailed arithmetic and proposed walk identification}
The following derivation uses the older partial-fraction notation so
that the complete-residue calculation can be stated precisely. The
normalization translation above makes it an argument about the same
three finite sums. The auxiliary rational function is used only for
this additional cancellation and the walk identification; it is not
needed for the coefficientwise denominator clearing in the main text.
\section{Exact objects and statement}
Let
\[
 F_n(t)=\frac{\prod_{u=0}^2(t-(10+u)n)_{10n}}
 {(t+n)_{16n+1}(t+2n)_{14n+1}},\qquad
 C_n=\frac{(15n)!^2}{(10n)!^3},\quad
 d_m=\mathop{\rm lcm}(1,\ldots,m).
\]
Write its partial fractions as
\[
 F_n(t)=\sum_{j=n}^{17n}\left(\frac{a_j}{(t+j)^2}+\frac{b_j}{t+j}\right),
 \qquad a_j=0\quad\hbox{outside }[2n,16n].
\]
Define, entirely by finite rational sums,
\begin{align}
 Q_n&=\sum_j a_j,\nonumber\\
 P_{3,n}&=\sum_j\left(a_jH^{(3)}_{j+2n}+\tfrac12 b_jH^{(2)}_{j+2n}\right),\label{arith-readout}\\
 P_{4,n}&=\sum_j\left(a_jH^{(4)}_{j+2n}+\tfrac13 b_jH^{(3)}_{j+2n}\right).\nonumber
\end{align}
Here $H_m^{(r)}=\sum_{k=1}^m k^{-r}$. Put
\[
 X_n=6C_nd_{19n}^4(Q_n,P_{3,n},P_{4,n}).
\]
This is the integer normalization whose gcd is estimated; it is not a
reduced denominator.

The eight factorial pairings used below are as follows. The column $b$
is a permutation of $(-2,-1,0,1,2)$ and realizes the displayed unordered
lengths $L=(10+b_0,11+b_1,12+b_2)$ and $D=(16-b_3,17-b_4)$.
\[
\begin{array}{c|c|c}
L&D&b\\\hline
(8,13,13)&(17,17)&(-2,2,1,-1,0)\\
(11,11,11)&(14,19)&(1,0,-1,2,-2)\\
(9,9,14)&(16,16)&(-1,-2,2,0,1)\\
(12,12,12)&(18,18)&(2,1,0,-2,-1)\\
(11,11,14)&(18,18)&(1,0,2,-2,-1)\\
(9,9,12)&(14,16)&(-1,-2,0,2,1)\\
(10,10,10)&(14,16)&(0,-1,-2,2,1)\\
(9,11,14)&(17,17)&(1,-2,2,-1,0)
\end{array}
\]
Let $\mathcal B$ denote these eight rows. Set
\begin{align*}
 c(x)&=2\lfloor15x\rfloor-3\lfloor10x\rfloor,\\
 \phi(x)&=\max\left(0,\max_{(L,D)\in\mathcal B}
 \left[c(x)-\sum_{d\in D}\lfloor dx\rfloor+
                  \sum_{l\in L}\lfloor lx\rfloor\right]\right).
\end{align*}
Call a prime eligible when
\[
 p>29,\qquad p^2>30n,\qquad p\nmid n.
\]
For eligible primes define
\[
 h_n(p)=\begin{cases}3&18n<p<19n,\\2&17n<p<18n,\\
 1&16n<p<17n,\\0&\text{otherwise},\end{cases}
 \qquad s_n(p)=\mathbf1_{8n<p<11n},
\]
and define the explicit positive integer
\begin{equation}\label{arith-G}
 G_n=\prod_{p\ {
m eligible}}p^{\phi(n/p)+h_n(p)+s_n(p)}.
\end{equation}
Only primes $p\leq19n$ occur with nonzero exponent.

\begin{theorem}\label{arith-main}
For every integer $n\geq1$, $X_n\in\Z^3$ and
\[
 G_n\mid\gcd(X_{n,0},X_{n,1},X_{n,2}).
\]
Furthermore
\[
 \liminf_{n\to\infty}\frac{\log G_n}{n}
 \geq29.0047746843482\ldots>29.004.
\]
For every fixed $A,B\in\Q$, reduce
\[
 \frac{AP_{3,n}+BP_{4,n}}{Q_n}=\frac{p_n}{q_n},\qquad q_n>0.
\]
Then, for all sufficiently large $n$,
\[
 q_n\leq\exp(112.5n).
\]
These assertions hold at every integer depth; no fractional lattice step
or restriction to a subsequence is used in this arithmetic argument.
\end{theorem}

\section{The finite sums, explained before estimating their gcd}
\subsection{The hypergeometric sequence is a finite sum}
The series in the question is
\[
 \mathcal F_n=\sum_{k=0}^{14n}
 \frac{(1+13n)_k(1+14n)_k(-14n)_k(-15n)_k(1+12n)_k}
 {(1+n)_k(1+2n)_k(1+3n)_k(1+4n)_k\,k!}.
\]
It terminates because $(-14n)_{14n+1}=0$. There is no convergence
question in this definition. Put
\[
 \rho_n=(-1)^n\frac{(12n)!(13n)!}
 {(2n)!(3n)!(4n)!n!(15n)!}.
\]
For $2n\leq j\leq16n$, direct cancellation of factorials gives
\begin{equation}\label{arith-explicit-a}
 a_j=(-1)^n
 \frac{(j+10n)!(j+11n)!(j+12n)!}
 {j!(j+n)!(j+2n)!(j-n)!(j-2n)!(16n-j)!(17n-j)!}.
\end{equation}
In particular $a_{2n}=\rho_n$. For $0\leq k<14n$ the ratio of successive
terms is
\[
 \frac{a_{2n+k+1}}{a_{2n+k}}=
 \frac{(1+12n+k)(1+13n+k)(1+14n+k)(k-14n)(k-15n)}
 {(1+k)(1+n+k)(1+2n+k)(1+3n+k)(1+4n+k)}.
\]
This is exactly the ratio of successive summands in $\mathcal F_n$.
Since their first terms also agree after multiplication by $\rho_n$,
\begin{equation}\label{arith-hypergeom-q}
 Q_n=\sum_{j=2n}^{16n}a_j=\rho_n\mathcal F_n.
\end{equation}

\subsection{Visible integrality: five binomial coefficients}
Multiplying \eqref{arith-explicit-a} by $C_n$ and grouping the factorials gives
the following identity, without an estimate or a divisibility conjecture:
\begin{equation}\label{arith-five-binomials}
 \boxed{\displaystyle
 C_na_j=(-1)^n
 \binom{j+10n}{10n}\binom{j+11n}{10n}\binom{j+12n}{10n}
 \binom{15n}{j-n}\binom{15n}{j-2n}.}
\end{equation}
Every binomial coefficient on the right has nonnegative integer arguments
in its allowed range. Therefore every summand is an integer. Consequently
\[
 C_nQ_n=(-1)^n\sum_{j=2n}^{16n}
 \binom{j+10n}{10n}\binom{j+11n}{10n}\binom{j+12n}{10n}
 \binom{15n}{j-n}\binom{15n}{j-2n}\in\Z.
\]
Equivalently $C_n\rho_n$ clears the denominator of the displayed
hypergeometric value $\mathcal F_n$. Notice the factor $\rho_n$:
$Q_n$ and $\mathcal F_n$ are not the same normalization.
For example, at $n=1,j=2$ the five factors are
$\binom{12}{10},\binom{13}{10},\binom{14}{10},\binom{15}{1},\binom{15}{0}$,
and their product has sign $-1$. This illustrates the same identity used
at every depth; no numerical example is needed for its proof.

\subsection{Why introduce the rational function $F_n(t)$?}
It is a compact way to encode the coefficients of three finite sums.
Its numerator has roots in the three intervals
$[1,10n]$, $[n+1,11n]$, $[2n+1,12n]$. Its denominator has roots in
$[-17n,-n]$ and $[-16n,-2n]$. A root appearing in both denominator
intervals gives a double pole. A root in just one gives a simple pole.
At a double pole, remove the factor $(t+j)^{-2}$:
\[
 A_j(t)=(t+j)^2F_n(t),\qquad a_j=A_j(-j),\qquad b_j=A'_j(-j).
\]
At a simple pole, $a_j=0$ and $b_j=((t+j)F_n(t))|_{t=-j}$.
These formulas define the partial-fraction coefficients by ordinary
polynomial evaluation and differentiation. Evaluating $A_j(-j)$ gives
\eqref{arith-explicit-a}.
If an entirely factorial-and-harmonic description is preferred, set
$H_m=\sum_{k=1}^m1/k$, with $H_0=0$. At a double pole,
\begin{align*}
 b_j=a_j\bigg(&\sum_{u=0}^2(H_{j+un}-H_{j+(10+u)n})
 +H_{j-n}-H_{17n-j}\\
 &+H_{j-2n}-H_{16n-j}\bigg),\qquad 2n\leq j\leq16n.
\end{align*}
This follows by logarithmically differentiating the factors left after
removing $(t+j)^2$. At the simple poles the corresponding formulas are
\[
 b_j=\begin{cases}
 \displaystyle(-1)^{j-n}
 \frac{\prod_{u=10}^{12}(j+un)!\,(2n-j-1)!}
 {j!(j+n)!(j+2n)!(j-n)!(17n-j)!(16n-j)!},&n\leq j<2n,\\[6pt]
 \displaystyle(-1)^{j-n+1}
 \frac{\prod_{u=10}^{12}(j+un)!\,(j-16n-1)!}
 {j!(j+n)!(j+2n)!(j-n)!(17n-j)!(j-2n)!},&16n<j\leq17n.
 \end{cases}
\]
Thus \eqref{arith-readout} specifies all three sequences using only finite
factorials and harmonic sums. The rational-function notation abbreviates
these explicit formulas; it does not supply an unproved integrality claim.

The denominator degree exceeds the numerator degree by two, so
$F_n(t)=O(t^{-2})$. Comparing the coefficient of $1/t$ in its partial
fractions gives $\sum_jb_j=0$. Now differentiate the partial fractions:
\[
 -\tfrac12 F'_n(t)=\sum_j\left(\frac{a_j}{(t+j)^3}
                         +\frac{b_j/2}{(t+j)^2}\right),\quad
 \tfrac16 F''_n(t)=\sum_j\left(\frac{a_j}{(t+j)^4}
                         +\frac{b_j/3}{(t+j)^3}\right).
\]
For $r>1$,
\[
 \sum_{k=2n+1}^{\infty}\frac1{(k+j)^r}
 =\zeta(r)-H^{(r)}_{j+2n}.
\]
There are only finitely many $j$, and each series converges absolutely.
Summing and using $\sum b_j=0$ proves the exact identities
\begin{equation}\label{arith-derivative-tails}
 \boxed{\displaystyle
 Q_n\zeta(3)-P_{3,n}=-\tfrac12\sum_{k=2n+1}^{\infty}F'_n(k),\qquad
 Q_n\zeta(4)-P_{4,n}=\tfrac16\sum_{k=2n+1}^{\infty}F''_n(k).}
\end{equation}
This explains both the zeta values and the harmonic corrections in
\eqref{arith-readout}. It does not yet estimate the size of these errors.

\subsection{Clearing the simple-pole coefficients, explicitly}
Define
\[
 P(t)=\prod_{u=0}^2\binom{t-un-1}{10n},\qquad
 L_c(t)=\frac{(15n)!}{\prod_{k=c}^{c+15n}(t+k)}.
\]
Regrouping precisely the same factors in $F_n$ gives
\begin{equation}\label{arith-reference-factorization}
 C_nF_n(t)=P(t)L_n(t)L_{2n}(t).
\end{equation}
For integer arguments, polynomial binomial coefficients are integers,
including negative arguments, because
$\binom{-a}{m}=(-1)^m\binom{a+m-1}{m}$ for $a\geq1$.
The simple partial fractions of $L_c$ are
\[
 L_c(t)=\sum_{j=c}^{c+15n}\frac{r_{c,j}}{t+j},\qquad
 r_{c,j}=(-1)^{j-c}\binom{15n}{j-c}\in\Z.
\]
Indeed the product of the other denominator factors at $t=-j$ is
$(-1)^{j-c}(j-c)!(c+15n-j)!$.
At a pole of $L_c$ let
\[
 s_{c,j}=\sum_{\substack{c\leq k\leq c+15n\\k\ne j}}
             \frac{r_{c,k}}{k-j}.
\]
Thus $L_c(t)=r_{c,j}/(t+j)+s_{c,j}+O(t+j)$ near $-j$.
At a shared pole, multiply these expansions and the Taylor expansion
$P(t)=P(-j)+(t+j)P'(-j)+O((t+j)^2)$. The result is
\begin{align*}
 C_na_j&=P(-j)r_{n,j}r_{2n,j},\\
 C_nb_j&=P'(-j)r_{n,j}r_{2n,j}
       +P(-j)(r_{n,j}s_{2n,j}+r_{2n,j}s_{n,j}).
\end{align*}
At a pole of only $L_c$, with the other factor denoted $L_{c'}$, it gives
\[
 C_nb_j=P(-j)r_{c,j}L_{c'}(-j),\qquad
 L_{c'}(-j)=\sum_{k=c'}^{c'+15n}\frac{r_{c',k}}{k-j}.
\]
Every nonzero difference $k-j$ here has absolute value at most $16n$.
Therefore $d_{16n}$ clears every $s_{c,j}$ and every displayed value
$L_{c'}(-j)$.

To treat the derivative, Vandermonde's polynomial identity says
\[
 \binom{k+x}{m}=\sum_{h=0}^{m}\binom{k}{m-h}\binom{x}{h}.
\]
The coefficient of $x$ in $\binom{x}{h}$ is $(-1)^{h-1}/h$ for $h\geq1$.
Differentiating at $x=0$ therefore gives
\[
 \left.\frac{d}{dx}\binom{x}{m}\right|_{x=k}
 =\sum_{h=1}^{m}\binom{k}{m-h}\frac{(-1)^{h-1}}h.
\]
Since each factor of $P$ has degree $10n$, the product rule shows that
$d_{10n}P'(-j)$ is an integer. As $d_{10n}\mid d_{16n}$, the preceding
explicit formulas prove
\[
 C_na_j\in\Z,\qquad d_{16n}C_nb_j\in\Z.
\]

\subsection{Clearing the two numerator sums term by term}
Write $D=d_{19n}$, $A_j=C_na_j$ and $B_j=DC_nb_j$.
Both $A_j$ and $B_j$ are integers. For $m\leq19n$,
\[
 D^rH_m^{(r)}=\sum_{k=1}^m(D/k)^r\in\Z,
\]
because every $k$ divides $D$. The exact normalized numerator formulas are
\begin{align*}
 6C_nD^4P_{3,n}
 &=\sum_j\left(6A_jD^4H^{(3)}_{j+2n}
                         +3B_jD^3H^{(2)}_{j+2n}\right),\\
 6C_nD^4P_{4,n}
 &=\sum_j\left(6A_jD^4H^{(4)}_{j+2n}
                         +2B_jD^3H^{(3)}_{j+2n}\right).
\end{align*}
Every term on both right-hand sides is an integer. Also
$6C_nD^4Q_n=6D^4\sum_jA_j\in\Z$.
This proves $X_n\in\Z^3$ with every denominator accounted for.
In particular, we have not assumed $C_nP_{3,n}$ or $C_nP_{4,n}$ is integral.

\section{Integrality from factorial blocks}
For any permutation $b$ of $(-2,-1,0,1,2)$, put
\[
 C_b(n)=\frac{\prod_{d\in D}(dn)!}{\prod_{l\in L}(ln)!}.
\]
The reference $C_n$ corresponds to $b=(0,-1,-2,1,2)$.

\begin{lemma}\label{arith-integral}
For each such pairing,
\[
 C_ba_j\in\Z,\qquad d_{19n}C_bb_j\in\Z,
 \qquad6C_bd_{19n}^4(Q_n,P_{3,n},P_{4,n})\in\Z^3.
\]
For the reference pairing the second assertion improves to
$d_{16n}C_nb_j\in\Z$.
\end{lemma}
\begin{proof}
The gamma quotient for $F_n$ has numerator shifts $-2n,-n,0,n,2n$
and denominator shifts $-10n,-11n,-12n,16n+1,17n+1$.
Explicitly, using $(x)_m=\Gamma(x+m)/\Gamma(x)$ gives
\[
 F_n(t)=\frac{\Gamma(t-2n)\Gamma(t-n)\Gamma(t)
                    \Gamma(t+n)\Gamma(t+2n)}
 {\Gamma(t-10n)\Gamma(t-11n)\Gamma(t-12n)
                    \Gamma(t+16n+1)\Gamma(t+17n+1)}.
\]
This identity is an identity of rational functions after cancellation;
we never evaluate the individual gamma factors at their poles.
Permuting the five numerator factors leaves the rational function
unchanged. Pair the first three with the first three denominator factors,
and the last two with the last two denominator factors.
Pairing these factors gives exactly
\[
 C_bF_n(t)=
 \prod_{i=0}^2\binom{t+b_in-1}{(10+i+b_i)n}
 \frac{((16-b_3)n)!}{\prod_{j=b_3n}^{16n}(t+j)}
 \frac{((17-b_4)n)!}{\prod_{j=b_4n}^{17n}(t+j)}.
\]
Each reciprocal factor has integral simple-pole residues, since
\[
 \mathop{\rm Res}_{t=-j}
 \frac{m!}{\prod_{k=a}^{a+m}(t+k)}
 =(-1)^{j-a}\binom m{j-a}.
\]
The polynomial factors are integer-valued on all integers. Their
individual degrees are at most $14n$, and their derivatives at integers
are cleared by $d_{14n}$. Indeed Vandermonde's polynomial identity gives
\[
 \left.\frac{d}{dx}\binom{x}{m}\right|_{x=k}
 =\sum_{h=1}^m\binom{k}{m-h}\frac{(-1)^{h-1}}h.
\]
The product rule gives the same derivative assertion for the polynomial
product, without replacing $14n$ by its total degree.

At a shared pole the double coefficient is a product of integer values
and residues. The simple coefficient adds derivative terms and terms
with denominators equal to differences of pole indices. At a pole of
only one reciprocal factor, only the latter denominators occur. All
index differences are at most $19n$. This proves the two coefficient
assertions; apparent extra poles cancel identically.
For the reference pairing, the two reciprocal intervals are
$[n,16n]$ and $[2n,17n]$, so differences are at most $16n$ and the
polynomial degrees are $10n$.

Finally $j+2n\leq19n$ at actual poles, and $d_{19n}^rH_m^{(r)}$ is
integral for $m\leq19n$. Formula \eqref{arith-readout} and the factor $6$
give the asserted integer triples.
\end{proof}

For eligible primes every factorial argument in $C_b$ is smaller than
$p^2$. Legendre's formula therefore gives
\[
 \vp(C_n/C_b)=c(n/p)-\sum_{d\in D}\lfloor dn/p\rfloor
                         +\sum_{l\in L}\lfloor ln/p\rfloor.
\]
Here $v_p(r)$ means the exponent of $p$ in a nonzero rational number $r$:
write $r=p^eu/v$ with $p$ dividing neither $u$ nor $v$, and set $v_p(r)=e$.
In particular an integer has nonnegative valuation; set $v_p(0)=+\infty$.
To see the divisibility consequence without skipping a step, define
\[
 Y_b=6C_b d_{19n}^4(Q_n,P_{3,n},P_{4,n})\in\Z^3.
\]
We have the exact vector equality $X_n=(C_n/C_b)Y_b$. Consequently,
for each coordinate $i$,
\[
 v_p(X_{n,i})=v_p(C_n/C_b)+v_p(Y_{b,i})\geq v_p(C_n/C_b).
\]
This does not require $C_n/C_b$ to be an integer: only its exponent of
the particular prime under consideration is relevant. For example,
if that exponent is two, all three entries of $X_n$ are divisible by $p^2$.
We already know $X_n$ is integral, so zero is another lower bound.
Taking the strongest bound separately at each prime proves
\begin{equation}\label{arith-pairbound}
 \min_i\vp(X_{n,i})\geq\phi(n/p).
\end{equation}
This is the factorial-normalization form of the arithmetic group method
of Rhin--Viola and Zudilin; no larger transformation group is assumed.
The freedom used here is simply to rewrite the same rational function
in eight ways. We neither divide the trajectory by four nor change its
integer vertices.

\section{The upper harmonic ranges}
We prove
\begin{equation}\label{arith-upperrange}
 \min_i\vp(X_{n,i})\geq h_n(p)\qquad(16n<p<19n).
\end{equation}
Here $C_n$ is a $p$-adic unit and $\vp(d_{19n})=1$.
If $j+2n<p$, Lemma~\ref{arith-integral} with $d_{16n}$ makes the corresponding
contribution $p$-integral before multiplying by $d_{19n}^4$.
For the remaining indices, $j>14n$.
At double poles,
\begin{equation}\label{arith-aformula}
 a_j=(-1)^n
 \frac{(j+10n)!(j+11n)!(j+12n)!}
 {j!(j+n)!(j+2n)!(j-n)!(j-2n)!(16n-j)!(17n-j)!}.
\end{equation}
At a right simple pole, the analogous formula for $|b_j|$ replaces
the factor $(16n-j)!$ in the denominator by $(j-16n-1)!$ in the
numerator. These formulas follow directly by removing the pole factors.
All three large numerator factorials contribute one $p$ each.
The only denominator factorials that can contain $p$ are
$j!,(j+n)!,(j+2n)!$. Logarithmic differentiation loses at most one
$p$ in passing from $a_j$ to $b_j$ at a double pole. Consequently:
\[
\begin{array}{c|c|c|c}
p\text{ range}&\vp(a_j)\text{ at active double poles}
 &\vp(b_j)\text{ there}&\vp(b_j)\text{ at active simple poles}\\\hline
(18n,19n)&\text{none}&\text{none}&\geq2\\
(17n,18n)&\geq2&\geq1&\geq1\\
(16n,17n)&\geq1&\geq0&\geq0
\end{array}
\]
An order-$r$ harmonic number loses at most $r$ powers of $p$ here.
Substitution into \eqref{arith-readout}, including the $d_{19n}^4$ factor,
proves \eqref{arith-upperrange} for $P_3,P_4$. For $Q$ the reference coefficient
integrality already gives valuation at least $4$ in $X_n$.

\section{One complete residue block}
\begin{lemma}\label{arith-block}
For an eligible prime $8n<p<11n$,
\[
 \vp(Q_n)\geq1,\qquad \vp(P_{3,n})\geq-1,\qquad
 \vp(P_{4,n})\geq-2.
\]
\end{lemma}
\begin{proof}
Write $F_n=\mathcal N/\mathcal D$ with monic polynomials. Numerator roots
are the three integer intervals $[1,10n],[n+1,11n],[2n+1,12n]$;
denominator roots are $[-17n,-n],[-16n,-2n]$.
For $0\leq k<p$ let $c_k,d_k$ count their roots congruent to $k$,
with multiplicity. We first establish
\begin{equation}\label{arith-support}
 c_k-d_k\geq-1,
 \qquad c_k-d_k=-1\ \Longleftrightarrow\
 k>12n-p,\quad k\geq2p-16n,\quad k\leq p-2n.
\end{equation}
For $0<k<p$ the exact counts are
\begin{align*}
 c_k={}&\mathbf1_{k\leq10n}+\mathbf1_{k>n}+\mathbf1_{k>2n}
       +\sum_{a=10}^{12}\mathbf1_{k\leq an-p},\\
 d_k={}&\mathbf1_{k\leq p-n}+\mathbf1_{k\leq p-2n}
       +\mathbf1_{k\geq2p-17n}+\mathbf1_{k\geq2p-16n}
       +\mathbf1_{k\geq3p-17n}.
\end{align*}
On $2n<k\leq p-2n$ the base numerator count is three and the denominator
count is at most four. It is four precisely when $k\geq2p-16n$;
there is no extra numerator root precisely when $k>12n-p$.
On $p-2n<k\leq p-n$ both counts equal three.
On $p-n<k<p$, $d_k=2+\mathbf1_{k\geq3p-17n}$ and
$c_k=2+\mathbf1_{k\leq10n}$; an extra denominator root forces
$p<17n/2$, hence $k<10n$.

For $0<k\leq n$, $c_k\geq2$. If $d_k\geq3$, then $p\leq9n$,
which makes $c_k=4\geq d_k$. For $n<k\leq2n$, $c_k\geq2$;
$d_k\geq3$ forces $p\leq19n/2$ and a third numerator root, while
$d_k=4$ forces $p\leq9n$ and a fourth numerator root.
Finally, at $k=0$, $c_0=3$ for $p<10n$ and $c_0=2$ for $p>10n$,
whereas $d_0=2+\mathbf1_{2p\leq17n}$. This proves \eqref{arith-support}.
On its support, $k>2n$, the numerator roots in this residue class are
$k,k,k$ and the denominator roots are $k-p,k-p,k-2p,k-2p$.

Reduction modulo $p$ now gives the polynomial
\[
 T(X)=(X^p-X)\frac{\overline{\mathcal N}(X)}
 {\overline{\mathcal D}(X)}\in\Fp[X],\qquad \deg T=p-2.
\]
This is a polynomial because each denominator multiplicity exceeds the
numerator multiplicity by at most one. Away from the support in
\eqref{arith-support}, $T(k)=0$.
More explicitly, $X^p-X=\prod_{k\in\Fp}(X-k)$ supplies exactly one
additional zero at every residue. It therefore cancels any one excess
denominator factor. Since $\deg\mathcal N=30n$ and
$\deg\mathcal D=30n+2$, the resulting degree is $p+30n-(30n+2)=p-2$.
On that support, separating the factors divisible by $p$ and using
$(X^p-X)'=-1$ yields the local identity
\begin{equation}\label{arith-local}
 pF_n(k+pV)\equiv-T(k)\frac{V^3}{(V+1)^2(V+2)^2}\pmod p.
\end{equation}
Outside the support this reduction is zero. The integral high-digit
pole locations in each residue class have pairwise differences of
absolute value at most two, hence are distinct modulo $p$. Partial
fraction extraction in this local identity is therefore legitimate.
It shows that $a_j/p$ and $b_j$ are $p$-integral, with both reducing
to zero outside the support. Notice that the coefficients of
$(V+h)^{-2}$ and $(V+h)^{-1}$ in $pF_n(k+pV)$ are $a_j/p$ and $b_j$.
For clarity, the three numerator roots $k,k,k$ contribute $(pV)^3$,
and the four denominator roots contribute
$p^4(V+1)^2(V+2)^2$. The prefactor $p$ cancels this remaining denominator
$p$. All other factors are units modulo $p$ and become constants.
Their quotient is $-T(k)$ because the missing factor from $X^p-X$
has derivative $-1$ at $k$. This proves \eqref{arith-local} directly.
Outside the support the numerator has at least as many local factors
as the denominator, and the prefactor $p$ leaves zero modulo $p$.

The elementary identity
\[
 \frac{V^3}{(V+1)^2(V+2)^2}
 =-\frac1{(V+1)^2}+\frac5{V+1}
  -\frac8{(V+2)^2}-\frac4{V+2}
\]
finishes the local calculation. The nodes $j=p-k$ and $j=2p-k$ have
harmonic indices $j+2n$ respectively less than $p$ and strictly between
$p$ and $2p$. Thus $p^rH_{j+2n}^{(r)}$ reduces respectively to $0$ and
$1$: all terms with $p\nmid k$ vanish after multiplication by $p^r$,
and the only possible remaining summand is $k=p$, which contributes one.
For the second node the displayed partial fractions give
\[
 a_{2p-k}/p\equiv8T(k),\qquad b_{2p-k}\equiv4T(k)\pmod p.
\]
Its contribution to $p^2P_{3,n}$ is therefore
$8T(k)+\tfrac12\,4T(k)=10T(k)$; its contribution to $p^3P_{4,n}$
is $8T(k)+\tfrac13\,4T(k)=\tfrac{28}3T(k)$.
The first node contributes zero because its harmonic index is less than
$p$. Classes outside the support contribute zero by \eqref{arith-local}.
This proves
\[
 p^2P_{3,n}\equiv10\sum_{k\in\Fp}T(k),\qquad
 p^3P_{4,n}\equiv\tfrac{28}{3}\sum_{k\in\Fp}T(k)\pmod p.
\]
For every polynomial of degree less than $p-1$, its sum over $\Fp$
is zero: check monomials, using the cyclic multiplicative group for
positive degrees and $\sum1=p=0$ for degree zero.
Since $\deg T=p-2$, both congruences vanish. Finally $Q_n=\sum a_j$
belongs to $p\Z_{(p)}$. This proves all three assertions.
Here $\Z_{(p)}$ denotes rational numbers whose reduced denominator is
not divisible by $p$. Thus $p^2P_{3,n}\equiv0\pmod p$ means exactly
$v_p(P_{3,n})\geq-1$, and $p^3P_{4,n}\equiv0\pmod p$ means
$v_p(P_{4,n})\geq-2$. This is the one step using cancellation among
different summands, rather than term-by-term denominator clearing.
\end{proof}

\section{Combining the divisors and evaluating their growth}
The explicit floor formula gives
\begin{equation}\label{arith-profilefacts}
0\leq\phi\leq3,\quad \phi(x+1)=\phi(x),\quad
\phi(x)=0\ (0<x<1/16),\quad
\phi(x)\leq c(x)+1\ (1/11<x<1/8).
\end{equation}
These are finite floor-function statements; the full exact partition is
given below. Periodicity follows also from $\sum D=\sum L$ in each row.
On the harmonic ranges $\phi(n/p)=0$, so \eqref{arith-upperrange} adds exactly
the stated $h_n(p)$. On $8n<p<11n$, Lemma~\ref{arith-block} gives the common
valuation of $X_n$ at least $c(n/p)+2\geq\phi(n/p)+1$.
The two additional prime ranges are disjoint. Elsewhere use
\eqref{arith-pairbound}. This proves divisibility by \eqref{arith-G}, without adding
two overlapping lower bounds incorrectly.

On one period, the nonzero intervals of $\phi$ are listed in the table;
on all omitted intervals it is zero. Endpoint values are irrelevant
for eligible primes: $n/p$ cannot have fractional part with reduced
denominator at most $19$.
\begin{longtable}{ccc@{\qquad}ccc}
\toprule $a$&$b$&$\phi$&$a$&$b$&$\phi$\\\midrule
1/15&1/13&1&1/13&1/12&2\\
1/12&1/10&3&2/15&2/13&1\\
2/13&3/17&2&3/17&2/11&1\\
2/11&1/5&2&1/5&4/15&1\\
4/15&5/17&3&5/17&3/10&2\\
1/3&2/5&2&2/5&4/9&1\\
5/11&7/15&1&7/15&9/19&3\\
9/19&1/2&2&8/15&6/11&1\\
6/11&3/5&2&3/5&2/3&1\\
2/3&11/16&3&11/16&9/13&2\\
9/13&7/10&3&11/15&14/19&2\\
14/19&3/4&1&3/4&4/5&2\\
4/5&13/16&1&9/11&16/19&1\\
13/15&7/8&2&7/8&15/17&3\\
15/17&9/10&2&14/15&17/18&2\\
17/18&18/19&1&&&\\\bottomrule
\end{longtable}
To verify the table, subdivide $[0,1]$ at every $i/j$, $1\leq j\leq19$.
Every floor in the definition is constant on each open interval, so one
exact rational midpoint evaluation per interval proves the entire table.

Let the sum below run over its nonzero intervals, with value $v$.
Exact rational addition gives
\begin{align*}
 S_8&=\sum_{(a,b,v)}v\sum_{k=0}^7
       \left(\frac1{k+a}-\frac1{k+b}\right)\\
 &=\frac{22763660987060292807135900464475002389291589797931744263}
 {1137911390967609212200766394165505379897899492140816000}
 >20.004.
\end{align*}
For each of these finitely many intervals, the prime number theorem
$\sum_{p\leq x}\log p=x+o(x)$ gives the contribution
$v((k+a)^{-1}-(k+b)^{-1})n+o(n)$ from
$n/(k+b)<p<n/(k+a)$. All omitted periods have nonnegative contribution.
The harmonic ranges contribute $(3+2+1)n+o(n)=6n+o(n)$, and the
residue block contributes $(11-8)n+o(n)=3n+o(n)$.
Excluding primes at most $\sqrt{30n}$ costs $O(\sqrt n\log n)=o(n)$,
because the exponents are bounded. Excluding prime divisors of $n$ costs
$O(\log n)$, and the finitely many primes at most $29$ cost $O(1)$.
Therefore
\[
 \liminf\frac{\log G_n}{n}\geq S_8+9>29.004.
\]
This proves the gcd-growth assertion for all depths, rather than
extrapolating from computed gcds.

\section{The reduced-denominator bound}
All nonzero $a_j$ in \eqref{arith-aformula} have the same sign. Hence $Q_n\ne0$.
For integer $A,B$, the three entries of $X_n/G_n$ are integers, so
\begin{equation}\label{arith-qbound}
 q_n\leq\frac{6C_nd_{19n}^4|Q_n|}{G_n}.
\end{equation}
For rational $A,B$, clearing their fixed denominators introduces only a
fixed factor on the right.

We give a separate uniform height bound. Set $f(x)=x\log x$, with $f(0)=0$,
and, for $2\leq y\leq16$, define
\[
 g(y)=\sum_{a=10}^{12}f(y+a)
 -f(y)-f(y+1)-f(y+2)-f(y-1)-f(y-2)-f(16-y)-f(17-y).
\]
Uniform Stirling bounds in \eqref{arith-aformula}, followed by summation over
at most $14n+1$ indices, give
\[
 \limsup\frac{\log(C_n|Q_n|)}n
 \leq30\log(3/2)+\max_{[2,16]}g.
\]
All $n\log n$ and linear Stirling terms cancel; the remaining error is
$O(\log n)$ uniformly, also at factorial argument zero.
On $(2,16)$, $g''<0$, since the three positive reciprocal terms
$1/(y+10),1/(y+11),1/(y+12)$ are already dominated by
$1/y,1/(y+1),1/(y+2)$ among the negative terms. Moreover
\[
 g'(y)=\log\frac{(y+10)(y+11)(y+12)(16-y)(17-y)}
 {y(y+1)(y+2)(y-1)(y-2)}.
\]
The difference between the numerator and denominator here is
$-30(15y^3+55y^2-1830y-11968)$.
The unique maximum lies between $a=12.014$ and $b=12.015$, as substitution
in this cubic verifies. For $y_0=(a+b)/2$, concavity gives
\[
 \max g\leq g(y_0)+\tfrac{b-a}{2}|g'(y_0)|.
\]
Rational logarithm bounds give
\[
 30\log(3/2)+g(y_0)+\tfrac{b-a}{2}|g'(y_0)|
 <65.451537<65.46.
\]
For completeness, every logarithm in this check is bounded after dyadic
scaling by the elementary convergent series, for $1\leq z\leq2$ and
$u=(z-1)/(z+1)$,
\[
 2\sum_{j=0}^{M-1}\frac{u^{2j+1}}{2j+1}
 \leq\log z\leq
 2\sum_{j=0}^{M-1}\frac{u^{2j+1}}{2j+1}
 +\frac{2u^{2M+1}}{(2M+1)(1-u^2)}.
\]
Taking $M=35$ proves the strict comparison by rational arithmetic.
Finally $\log d_{19n}=19n+o(n)$, so \eqref{arith-qbound} implies
\[
 \limsup\frac{\log q_n}n
 \leq65.46+76-(S_8+9)
 =112.45522531565\ldots<112.5.
\]
This finishes the proof of Theorem~\ref{arith-main}.

\section{Relation to the requested forward walk}\label{arith-walk-section}
Let $H$ be the user's native $4\times4$ ${}_5F_4$ trajectory matrix and
$W_n^{(5)}=H(0)\cdots H(n-1)$, for
\[
 w_0=(1,1,0,0,1;1,1,1,1),\qquad
 v=(13,14,-14,-15,12;1,2,3,4).
\]
Define
\[
 \kappa_n=\frac{(12n)!(13n)!}
 {(2n)!(3n)!(4n)!n!(15n)!},\qquad \rho_n=(-1)^n\kappa_n.
\]
The exact readout identity is
\begin{equation}\label{arith-bridge}
 \begin{pmatrix}Q_n\\P_{3,n}\\P_{4,n}\end{pmatrix}
 =\rho_n
 \begin{pmatrix}1&0&0&0\\0&0&-1/2&1/4\\0&0&0&-1/6\end{pmatrix}
 W_n^{(5)}e_0.
\end{equation}
In particular the two-row expression in the question is
$(AP_{3,n}+BP_{4,n})/Q_n$ exactly. The integer normalization of the
three user's rows $(12,0,0,0),(0,0,-6,3),(0,0,0,-2)$ is
$\rho_n C_nd_{19n}^4/2$, producing $X_n$.
Thus the gcd estimate is for the actual rational readouts, with their
scalar normalization fully specified.

Here is a derivation of the identities used in \eqref{arith-bridge}. In the
equivalent Meijer--G coordinates the start is
$(1,1,0,0,0;0,0,0,0,0)$ and the direction is
$(16,17,-11,-12,-10;2,1,0,-1,-2)$.
At an integer vertex use the rational gamma quotient
\[
 R_w(t)=(-1)^{a_0+a_1}\frac{\prod_j\Gamma(t+b_j)}{\prod_i\Gamma(t+a_i)},
 \qquad J_w(t)=R_w(t)\prod_j(t+b_j).
\]
The gamma functional equation gives the elementary contiguous
multipliers, with the jet variable represented by $-t$, and
\[
 R_w(t)\prod_i(t+a_i-1)=J_w(t-1).
\]
Hence the degree-four companion relation is the translate difference
$J_w(t-1)-J_w(t)$.
Use the coordinate-block path: lower $a_2,a_3,a_4$, raise $a_0,a_1$,
raise $b_0,b_1$, then lower $b_3,b_4$. The leading coefficient
$5-\sum a_i+\sum b_j$ stays at least $3$. For the inverse $a$ moves,
the old positive $a_i$ never equals a $b_j=0$; for the inverse $b$ moves,
the new $b_j\in[-2n,-1]$ never equals an $a_i-1$.
All required inverse generators therefore exist.

Use the two periodic kernels $K_2(t)=\pi^2/\sin^2\pi t$ and
$K_3(t)=\pi^3\cos\pi t/\sin^3\pi t$.
During the first three blocks take $0<\Re t<1$.
At their end, the rational numerator has triple zeros at
$1,\ldots,10n$, so move the contour to $2n<\Re t<2n+1$ before the last
block. No residue is crossed. To annihilate the translate difference,
$J_w$ has a triple zero at $0$ on the first contour and at $2n$ on the
last contour. At the last endpoint a lost zero of $R_w$ at $2n$ is
restored by the factor $t+b_4$ in $J_w$. Horizontal integrals vanish
by exponential decay of the kernels. Thus the period functionals obey
the contiguous matrices and their companion reductions along this path.
The sums of double-pole and simple-pole coefficients also annihilate
integer translate differences and polynomials, and obey the same matrices.

Initially $R_w=t^{-2}$. On the jet basis $(-t)^rR_w$, $0\leq r\leq3$,
the four functional rows are
\[
 (1,0,0,0),\quad(0,-1,0,0),\quad
 (2\zeta(3),-\zeta(2),1,-1/2),\quad
 (-3\zeta(4),\zeta(3),0,-1/2).
\]
The first two are the coefficient sums. The period values follow by
closing right and differentiating simple partial fractions; the last
entries also follow from $K_2=-\pi(\cot\pi t)'$, reflection about $1/2$,
and $K_3=-K_2'/2$.
At the endpoint $R_w=(-1)^nF_n$, and the periods are
$(-1)^n2(Q_n\zeta(3)-P_{3,n})$ and
$(-1)^n[-3(Q_n\zeta(4)-P_{4,n})]$.
The simple coefficient sum is zero because $F_n=O(t^{-2})$.
Comparing these functional values gives the Meijer--G version of
\eqref{arith-bridge}, without using independence of any zeta values.

Finally the rational gauge between the two contiguous systems is
\[
 M_G(n)=\frac{\kappa_{n+1}}{\kappa_n}K(n)^{-1}H(n)K(n+1),
 \qquad K(n)_{ij}=\begin{cases}\binom ji(2n)^{j-i}&i\leq j,\\0&i>j.
 \end{cases}
\]
It follows by the parameter substitution
\[
 x=(1+2n-a_2,1+2n-a_3,1+2n-a_0,1+2n-a_1,1+2n-a_4),
 \quad y=(1+2n-b_1,1+2n-b_2,1+2n-b_3,1+2n-b_4)
\]
in the contiguous relations and the binomial change of jet origin.
The accompanying algebraic verifier also checks all sixteen entries
of this identity as rational functions, not merely at sample depths.
Since $K(0)=I$ and $K(n)e_0=e_0$, telescoping the gauge gives
$W_n^{(G)}e_0=\kappa_nW_n^{(5)}e_0$, which proves \eqref{arith-bridge}.

As a scalar check, \eqref{arith-aformula} has term ratio
\[
 \frac{a_{2n+k+1}}{a_{2n+k}}
 =\frac{(1+12n+k)(1+13n+k)(1+14n+k)(k-14n)(k-15n)}
 {(1+k)(1+n+k)(1+2n+k)(1+3n+k)(1+4n+k)}.
\]
Thus $Q_n=\rho_n\,{}_5F_4(1+13n,1+14n,-14n,-15n,1+12n;
1+n,1+2n,1+3n,1+4n;1)$, with exactly the family in the question.

\section{All rational initial rows of the companion}
The common divisor also applies to arbitrary fixed rational initial
rows, not just the two displayed numerator rows. Let $C(n)$ be the
three-dimensional companion and
\[
 U_0=[e_0,H(0)e_0,H(0)H(1)e_0],\qquad
 R=\begin{pmatrix}1&0&0&0\\0&0&-1/2&1/4\\0&0&0&-1/6\end{pmatrix},
 \quad S=RU_0.
\]
The companion identity gives $W_n^{(5)}e_0=U_0Z_ne_0$, where
$Z_n=C(0)\cdots C(n-1)$. The matrix $S$ is invertible over $\Q$.
A short exact certificate is that its entries are $101$-integral and
\[
 S\pmod {101}=\begin{pmatrix}1&89&15\\0&35&64\\0&98&61\end{pmatrix},
 \qquad \det S\equiv4\pmod {101}.
\]
These entries are obtained from the first two forward steps; the
verifier reconstructs those columns and compares the saved companion
coboundary exactly.

For any finite list of fixed rational initial rows $u_1,\ldots,u_k$,
choose a positive integer $d$ clearing all entries of $u_iS^{-1}$.
Then \eqref{arith-bridge} and Theorem~\ref{arith-main} imply
\[
 \boxed{\displaystyle
 \frac{6d\,C_nd_{19n}^4\rho_n}{G_n}\,u_iZ_ne_0\in\Z
 \qquad(1\leq i\leq k,\ n\geq1).}
\]
Indeed the left side equals $d\,u_iS^{-1}(X_n/G_n)$.
Thus the explicit divisor and its uniform density bound are inherited
by every finite collection of rational initial conditions after a
fixed denominator clearing. This does not presume linear independence
of their limiting ratios.


\begin{thebibliography}{9}
\bibitem{DLMF} NIST Digital Library of Mathematical Functions,
\href{https://dlmf.nist.gov/16.3.E1}{Eq.~16.3.1},
\href{https://dlmf.nist.gov/5.4.E14}{Eq.~5.4.14}, and
\href{https://dlmf.nist.gov/5.2}{\S5.2},
\href{https://dlmf.nist.gov/5.5.E3}{Eq.~5.5.3},
\href{https://dlmf.nist.gov/5.12.E1}{Eq.~5.12.1},
\href{https://dlmf.nist.gov/5.13.E1}{Eq.~5.13.1}, and
\href{https://dlmf.nist.gov/16.17.E1}{Eq.~16.17.1}.
\bibitem{RV} G.~Rhin and C.~Viola,
\emph{The group structure for $\zeta(3)$},
Acta Arith. \textbf{97} (2001), 269--293.
\url{https://www.math.ru.nl/~zudilin/zw/Rhin-Viola_AA_2001.pdf}.
\bibitem{Zudilin} W.~Zudilin,
\emph{Arithmetic of linear forms involving odd zeta values}.
\url{https://arxiv.org/abs/math/0206176}.
\end{thebibliography}
\end{document}
